\documentclass[11pt]{article}

\usepackage[T1]{fontenc}
\usepackage[margin=1in]{geometry}
% \usepackage[round]{natbib}
\usepackage[authoryear,round]{natbib}
\usepackage{amsmath,amssymb,amsthm}
\usepackage{mathtools}
\usepackage{bm}
\usepackage{graphicx}
\usepackage{booktabs}
\usepackage{subcaption}
\usepackage{algpseudocode}
\usepackage{algorithmicx}
\usepackage{algorithm}
\usepackage{float}
\usepackage{placeins}
\usepackage[hidelinks]{hyperref}
\usepackage{xcolor}

% ---- theorem environments ----
\theoremstyle{plain}
\newtheorem{theorem}{Theorem}
\newtheorem{proposition}{Proposition}
\newtheorem{corollary}{Corollary}
\newtheorem{lemma}{Lemma}
\theoremstyle{definition}
\newtheorem{definition}{Definition}
\newtheorem{assumption}{Assumption}
\newtheorem{remark}{Remark}
\newtheorem{hypothesis}{Hypothesis}

% ---- convenient macros ----
\newcommand{\R}{\mathbb{R}}
\newcommand{\E}{\mathbb{E}}
\newcommand{\thetab}{\bm{\theta}}
\newcommand{\wb}{\bm{w}}
\newcommand{\xb}{\bm{x}}
\newcommand{\fphys}{f_{\mathrm{phys}}}
\newcommand{\gnn}{g_{\mathrm{NN}}}
\newcommand{\Identif}{\mathcal{I}}
\newcommand{\Out}{\mathcal{Y}}
\newcommand{\Uclass}{\mathcal{U}}
\newcommand{\Feff}{\mathcal{F}_{\mathrm{eff}}}
\newcommand{\Gres}{\bm{G}}
\DeclareMathOperator{\col}{col}
\DeclareMathOperator{\spn}{span}
\DeclareMathOperator{\rank}{rank}
\newcommand{\todo}[1]{\textcolor{red}{[#1]}}

\newcommand{\method}{PlanGrad}
\newcommand{\bx}{\bm{x}}
\newcommand{\bu}{\bm{u}}
\newcommand{\bw}{\bm{w}}

\begin{document}

\title{\bf Does prediction accuracy buy separation?\\
	A closed-loop audit of prediction, planning, and manoeuvring
	authority in low-altitude corridors}

\author{}
\date{}

\maketitle

\begin{abstract}
	Trajectory predictors for low-altitude traffic are selected by
	open-loop displacement error, and coupling them to the planner
	during training is expected to improve safety further. Neither
	premise has been tested where separation is enforced as an explicit
	constraint. We develop \method, a differentiable closed-loop testbed
	that composes multi-agent prediction, control-barrier-constrained
	planning, and six-degree-of-freedom eVTOL dynamics, so that a single
	component can be exchanged while the encounter is replayed. Across
	$200$ paired conflict encounters, a twenty-five-fold spread in
	average displacement error ($0.83$--$20.90$\,m) yields conflict
	rates within $1.5$ percentage points ($11.0$--$12.5\%$), and
	displacement ranking does not carry over to closed-loop separation.
	Routing gradients through the planner produces a small, locally
	conservative shift in prediction error near closest approach
	($-0.164$\,m) with no detectable operational increment. Residual
	conflicts are not resolved by replacing the forecast with the true
	future trajectory; they instead track manoeuvring authority and
	available response time, disappearing once detection lead time
	reaches $3$\,s, invariantly across planning horizons of
	$1.6$--$5.0$\,s. Cross-evaluating predictors under a second planner
	shows that differences compressed to $2.5$ points under deployment
	widen to $42$ points elsewhere. Corridor simulations with mixed
	avoidance logics reduce delay, including for unequipped aircraft,
	while corridor-retention costs grow. Operational value is a property
	of the composed system, not of the predictor alone.
\end{abstract}

\noindent\textbf{Keywords:} urban air mobility; trajectory prediction;
control barrier functions; conflict detection and resolution;
decision-focused learning; mixed autonomy

\section{Introduction}

Low-altitude mobility is becoming a new layer of the urban
transportation network. Electric vertical take-off and landing vehicles,
cargo drones, and inspection unmanned aerial vehicles are adding
passenger, logistics, and service flights to already complex urban
systems. These operations bring new demand patterns, infrastructure
needs, capacity constraints, and operating rules~\citep{coppola2024uam,wang2023uam,pak2024uam}. The central challenge is no longer only whether an individual aircraft can fly autonomously. It is whether many aircraft can share limited urban airspace while maintaining separation and providing reliable service. In transportation systems, safety is typically shaped jointly by vehicle capabilities, infrastructure, operating rules, and automated behavior, rather than determined by any single component. For low-altitude traffic, however, the relative contribution of these layers is not yet established, and the answer matters. It determines where research effort, evaluation
practice, and eventually admission and airspace standards should
concentrate.

These factors interact most directly in the tactical management of
conflicts inside designated low-altitude corridors. Corridors create
shared operating regions in which aircraft traveling between vertiports
and other service nodes must resolve crossing and merging encounters
within constrained three-dimensional
space~\citep{tang2023strategic,wang2023quasidynamic}. The outcome of
such an encounter depends on several layers simultaneously. Onboard
autonomy must predict the surrounding aircraft and select an avoidance
maneuver, the vehicle must have sufficient maneuvering capability to
execute it, and airspace and surveillance design determine how much time
remains for the response. These encounters are further complicated by wind disturbances and imperfect observation of the surrounding traffic~\citep{nithya2024wind}. In an active corridor, stopping or hovering cannot be treated as a general recovery option. Many conflicts must therefore be resolved while both aircraft remain in motion, and the available response window is often measured in seconds rather than minutes~\citep{kuchar2000review,chen2024integrated}. Across repeated corridor operations, these vehicle-level decisions shape corridor-level loss-of-separation risk, delay, and usable capacity.

Across much of the current low-altitude autonomy literature, these
vehicle-level decisions are produced by a common modular structure.
A predictor estimates the future motion of surrounding aircraft, and
a planner uses those estimates to select a trajectory intended to
maintain separation. The two components are usually developed and evaluated separately. Prediction models are compared mainly using average displacement error (ADE) or final displacement error (FDE), usually without considering how their forecasts affect downstream control decisions~\citep{zhang2023flight,guo2024integrating}. Planning methods, in turn, are assessed using closed-loop measures such as conflict rate, minimum separation, and control effort~\citep{lindemann2023safe}. This evaluation practice rests on a common premise. More accurate forecasts are expected to support safer operation, and allowing the planner to guide predictor training is expected to provide further operational gains~\citep{huang2023differentiable,zhang2025bridging,mandi2024decisionfocused}.

That premise is plausible, but it has not been established for low-altitude traffic operations. Within a prediction-then-planning pipeline, the operational effect of a forecast depends on how it changes the actions selected by the planner. ADE averages displacement errors across the prediction horizon, whereas FDE reports the error at the final prediction step. What matters for closed-loop safety, however, is whether a prediction error changes a maneuver at a time and location relevant to separation. Neither metric reveals whether this occurs. A planner that enforces separation through explicit constraints may also respond similarly to forecasts with very different displacement errors. It therefore remains unclear whether improvements in open-loop prediction, or training the predictor through the planner, translate into additional benefits in closed-loop operation. This raises two questions. Does the ranking of predictors by displacement accuracy carry over to their ranking by closed-loop safety? Does training the predictor through a separation-constrained planner provide an operational benefit beyond ordinary fine-tuning on the same closed-loop training encounters?

Answering these questions requires a controlled closed-loop setting. Recorded flight tracks show what occurred under one sequence of actions, but they do not reveal how the same encounter would have unfolded under a different predictor or planner. We therefore develop \method, an interactive testbed that predicts the motion of surrounding aircraft, simulates the six-degree-of-freedom motion of the ego aircraft, and selects control actions subject to an explicit minimum-separation constraint. Prediction uncertainty is calibrated before it enters the planner, and the separation constraints are implemented using control barrier functions (CBFs). The testbed allows the same encounter to be replayed while changing one component at a time. It also includes a reference case in which the future trajectory of the neighbouring aircraft is known, helping distinguish conflicts caused by prediction error from those that remain even with perfect future information. Differentiability serves a narrower purpose. It allows the operational consequences of the planner's actions to be propagated back into predictor training. This enables a controlled comparison between planner-guided training and ordinary fine-tuning using the same encounter data and training conditions. We initially expected planner-guided training to provide an additional benefit in closed-loop operation.

The analysis follows the layers introduced above through a sequence of
controlled comparisons and robustness analyses. We first vary the
predictor while holding the planner fixed to assess the role of
prediction. We then retain the predictor and alter the planning layer
to examine the contribution of explicit separation constraints and
available manoeuvring authority. A reference case with known future
trajectories distinguishes conflicts that improved prediction could
prevent from those that persist even with perfect future information.
Plant--model mismatch and stronger wind disturbances test robustness
to departures from the nominal operating model. We additionally
cross-evaluate Stage~2 predictors under two planning configurations to
examine whether predictor-side differences retain the same closed-loop
expression when the downstream planner changes. A sweep of available
detection lead time then examines the value of additional response
time. Finally, corridor simulations extend the analysis from individual
encounters to mixed operations involving aircraft governed by different
avoidance logics.



Together, these analyses connect vehicle-level prediction and control to questions faced directly by traffic operators and airspace designers. They examine how prediction models should be evaluated when used for tactical separation and whether minimum maneuver-performance requirements should enter corridor admission criteria. They also consider how surveillance and airspace design should be assessed through the response time they provide, and what coordination challenges may arise when aircraft using different avoidance logics share a corridor.

The contributions of this paper are summarized as follows.

\begin{itemize}
\item[(i)] We formulate tactical conflict management in low-altitude
		corridors as a closed-loop transportation problem and develop
		\method, a differentiable testbed that integrates multi-agent
		trajectory prediction, separation-constrained planning, and
		six-degree-of-freedom eVTOL dynamics within a common closed loop.

	\item[(ii)] We design a controlled evaluation methodology for this
problem. Single-change comparisons, paired encounter-level statistics,
and a reference case with known future trajectories connect the
open-loop displacement metrics by which predictors are usually judged
to controller-in-the-loop separation outcomes. A matched training
comparison, complemented by a directional decomposition of prediction
errors, audits planner-guided predictor training against ordinary
fine-tuning on the same encounter data, while cross-configuration
evaluation tests whether predictor differences retain the same
operational expression when the downstream planner changes.

	\item[(iii)] We broaden the evaluation across operating conditions
		and scales, from plant--model mismatch, wind disturbances, and
		conflict-detection lead time to corridor operations in which
		aircraft governed by different avoidance logics interact, thereby
		connecting tactical conflict resolution to corridor-level safety,
		delay, and throughput.

\end{itemize}


% =====================================================================
% Section 2: Related Work
% =====================================================================
\section{Related work}
\label{sec:related}

This study draws on three research streams that have advanced largely
in parallel. Aircraft trajectory prediction supplies the models under
test, safe planning under predictive uncertainty supplies the planning
layer, and the joint training of prediction and planning, developed in
automated driving, supplies the hypothesis this study examines.
Reviewing them in turn also explains why their intersection remains
unexamined for low-altitude traffic operations.

\subsection{Trajectory prediction and the yardstick it inherited}
\label{sec:related-prediction}

Deep learning has reshaped aircraft trajectory prediction over the
past few years. In commercial aviation, Bayesian deep learning has
brought weather-induced uncertainty into the forecast
itself~\citep{pang2021data}, time--frequency wavelet
representations have improved prediction by capturing non-stationary
motion components~\citep{zhang2023flight}, and prediction models have
begun to absorb air traffic control voice instructions as an
additional input channel~\citep{guo2024integrating}. Trajectory prediction
for low-altitude traffic remains considerably less developed. Recent
UAM-specific studies have begun to narrow the gap, including
deep-learning frameworks that pair trajectory prediction with
downstream collision prediction~\citep{kim2025dlcollision} and
graph-based models that forecast drone states from Remote~ID data to
support separation monitoring~\citep{wen2025bidgcnllm}. Road traffic,
by contrast, already supports a mature interaction-aware multi-agent
prediction literature, in which heterogeneous road users are modelled
jointly through graph attention over kinematically consistent
representations~\citep{sheng2024kinematics} and long-horizon forecasts
are grounded in recurring vehicle-following behaviour
patterns~\citep{an2025longterm}. An equivalent body of work has yet to
form for low-altitude traffic.

Across learning-based trajectory prediction, performance is still
assessed primarily through open-loop geometric errors, including
average and final displacement error (ADE and FDE) in multi-agent
motion forecasting and RMSE-based measures in aircraft trajectory
prediction. These metrics quantify geometric accuracy, but they do
not reveal whether a prediction error changes a safety-critical
planning decision. The concern is beginning to surface in adjacent
communities: \citet{liu2025risknet} integrate trajectory prediction
with an explicit interaction-risk representation, illustrating a
broader shift from purely geometric forecasting toward operationally
relevant risk assessment. What is still missing is a direct test, for
low-altitude operations, of whether ranking predictors by
displacement accuracy also ranks them by closed-loop safety.
Providing that test requires a planner in the loop, which is where
the second stream enters.

\subsection{Safe planning that treats the predictor as given}
\label{sec:related-planning}

Research on planning around uncertain traffic has produced two
principled ways of consuming a forecast. Conformal prediction wraps
the forecast in distribution-free uncertainty sets and enforces
separation against the set rather than the point
estimate~\citep{lindemann2023safe}, with extensions to partially
observable settings~\citep{sheng2024safe} and to calibration under
distribution shift~\citep{huang2025cuqds}. Control barrier functions
instead encode separation through conditions that can render a
prescribed safe set forward invariant when the constraints remain
feasible. The construction was established for multi-robot collision
avoidance through real-time quadratic
programs~\citep{wang2017safety}, rests on a mature theory of safe
learning-based control~\citep{brunke2022safe}, and has been extended
to robust formulations coupled with reinforcement
learning~\citep{emam2022sablas}, to learned certificates for
black-box dynamics~\citep{qin2022sablas}, and, in transportation, to
platoon-level safety in mixed-autonomy road
traffic~\citep{zhou2024platoon}.

Low-altitude traffic management has meanwhile developed its own
planning stack. Conflict detection and resolution has a modeling
tradition reaching back to \citet{kuchar2000review}. Recent work
spans the strategic layer, where flight plans are deconflicted before
departure~\citep{tang2023strategic}, the network layer, where traffic
is assigned across routes to manage complexity and
congestion~\citep{wang2023quasidynamic}, and the design of the
airspace structures themselves~\citep{bauranov2021designing}. At the
tactical layer, computational guidance methods resolve conflicts
online during flight~\citep{yang2021autonomous}, including
formulations that scale to many agents while carrying explicit
separation assurance~\citep{yang2020scalable}, and integrated
architectures now couple strategic demand management with
learning-based tactical deconfliction~\citep{chen2024integrated}.

A common feature of much of the planning literature reviewed above is
that the predictor remains a fixed upstream module. Its outputs are
passed to the planner, but the planner does not influence how the
predictor is trained. The planning layer is engineered to tolerate
whatever error the predictor delivers, which raises, but does not
answer, the question of how much operational value a better predictor
then adds. The third stream supplies a candidate answer.

\subsection{Integrated prediction and planning, and the untested case}
\label{sec:related-integrated}

In automated driving, the separation between prediction and planning
has begun to dissolve. \citet{huang2023differentiable} differentiate
through an optimization-based planner so that planning cost helps
supervise the predictor, and report improved closed-loop planning and
prediction relative to a conventional two-stage pipeline. More
broadly, end-to-end architectures increasingly treat prediction and
planning within one trainable model~\citep{zhang2025bridging}. These
developments are closely related to decision-focused learning, which
is formalized in frameworks such as smart
predict-then-optimize~\citep{elmachtoub2022smart}, with theory
characterizing when two-stage learning forfeits performance and when
it is already
optimal~\citep{cameron2022perils,mandi2024decisionfocused}.
Differentiable optimization layers such as OptNet provide one route
for propagating decision losses through
optimization~\citep{amos2017optnet}, while software frameworks such
as PyEPO make predict-then-optimize training available across a
broad range of linear and integer programs~\citep{tang2022pyepo}.

A direct transfer of these gains to low-altitude corridors is not
straightforward. eVTOL flight involves six-degree-of-freedom
rigid-body dynamics with explicit actuator and flight-envelope
constraints~\citep{acheson2024guam,campbell2024pyguam}, and urban
airflow and weather introduce disturbances that vary strongly with
the built environment~\citep{nithya2024wind,reiche2021initial}, with
wind conditions already shown to constrain eVTOL operations and
vertiport traffic flows in city-specific case
studies~\citep{schweiger2023wind}.
Conflict resolution must also be carried out in free
three-dimensional space~\citep{rezaee2024collision}, with separation
maintained continuously against forecasts of surrounding traffic,
rather than within a planar lane structure. More importantly, the
differentiable driving framework most directly comparable to the
present study couples prediction and planning through a learned
planning cost rather than an explicit minimum-separation
constraint~\citep{huang2023differentiable}. Existing evidence
therefore does not establish whether planner-guided prediction
retains its value when separation is enforced as a constraint.

End-to-end learning also exists in the aerial domain, but it
addresses a different problem. Representative work maps onboard
sensing directly to collision-free trajectories for a single vehicle,
enabling agile flight through cluttered
environments~\citep{loquercio2021learning}; it does not forecast the
behavior of surrounding traffic, and it involves no separation
constraint against such traffic. We have not identified prior work
that trains a surrounding-traffic predictor through an explicit
separation-constrained planner for low-altitude operations, and the
closed-loop instrument of Section~\ref{sec:method} is built to study
exactly that.

One further tradition informs the system-level part of this study. Transportation research has long examined what happens when vehicles under different control logics share a facility, from the dependence of traffic flow stability and throughput on the penetration of connected and automated vehicles~\citep{talebpour2016influence} to field experiments in which a single controlled vehicle reshapes the surrounding
traffic~\citep{stern2018dissipation}. That literature has since
established that outcomes in mixed fleets depend not only on the share
of automated vehicles but on how they are spatially arranged, since
the tendency of automated vehicles to cluster alters the capacity
attained at a given penetration rate~\citep{zhao2024clustering}; and
reviews of continuum traffic-flow modelling in the connected and
automated era emphasise that aggregate models must be re-examined
rather than reused unchanged when a new class of controlled vehicles
enters the stream~\citep{mohammadian2023continuum}. Our mixed-equipage
corridor experiment poses the low-altitude form of that question,
without presuming that the ground-traffic answers carry over.

\section{Problem Formulation}
\label{sec:problem}
 
\subsection{Operational setting and scope}
\label{sec:problem-setting}
 
We consider a single aircraft under our control, referred to as the
\emph{ego aircraft}, flying through a designated low-altitude corridor
between two vertiports. The corridor supplies a reference route
$\bm{r}_{0:T_{\mathrm{end}}}$, a nominal path with desired passage times
that the aircraft is asked to follow. The state of the ego aircraft at
discrete time $t$ is denoted $\bx_t \in \R^{n_x}$ and collects its
position $\bm{p}_t \in \R^3$ together with its velocity and attitude. The
control input $\bu_t \in \mathcal{U} \subset \R^{n_u}$ comprises the
thrust and moment commands available to the vehicle, and the admissible
set $\mathcal{U}$ encodes its actuator limits. Local wind acts on the
vehicle as an unmeasured disturbance $\bw_t \in \R^3$; the decision
system does not observe the wind realization. The motion of the ego aircraft is governed by
discrete-time dynamics
\begin{equation}
	\label{eq:dynamics}
	\bx_{t+1} = g(\bx_t,\, \bu_t,\, \bw_t),
\end{equation}
which is the only property of the vehicle needed to state the problem;
the specific vehicle model and its numerical implementation are described
in Section~\ref{sec:experiments}.
 
The ego aircraft shares the corridor with other traffic. At each time $t$
it observes a set of neighbouring aircraft $\mathcal{N}_t$, and for each
neighbour $i \in \mathcal{N}_t$ it holds a window of recent observed
states $H^i_t$, obtained from broadcast position reports or onboard
tracking. The collective history is denoted
$H_t = \{H^i_t\}_{i \in \mathcal{N}_t}$. Future motion of the neighbours
is not known to the ego aircraft.
 
Operations in the corridor are evaluated against a prescribed minimum
separation threshold $d_{\mathrm{sep}} > 0$, a keep-clear distance that
the operation is expected to maintain around every other aircraft. At
every instant, the ego aircraft is expected to remain at least
$d_{\mathrm{sep}}$ away from each neighbouring aircraft,
\begin{equation}
	\label{eq:sep-requirement}
	\big\| \bm{p}_t - \bm{p}^i_t \big\|_2 \;\geq\; d_{\mathrm{sep}}
	\qquad \text{for all } t \text{ and all } i \in \mathcal{N}_t,
\end{equation}
where $\bm{p}^i_t$ is the true position of neighbour $i$. An encounter
in which the distance falls below this threshold is classified as a loss
of separation, even though no physical collision has occurred. The
requirement~\eqref{eq:sep-requirement} is stated on true positions; how a
system that only holds uncertain forecasts of those positions can act so
as to satisfy it is precisely the design problem of the planning layer.
 
Two remarks fix the scope of the study. First, joint planning among
several controlled aircraft is outside the present scope. The
controlled encounter-level experiments specialize to one neighbouring
aircraft so that prediction--planning interactions and feasibility can
be attributed in isolation. The formulation itself allows several
neighbours, and the later corridor-scale extension permits each
equipped aircraft to consider several neighbours simultaneously, but
that extension is exploratory and does not perform the controlled
multi-constraint feasibility attribution used in the single-neighbour
experiments. Second, the static environment enters only through the
corridor geometry and the reference route; richer map information,
such as building shapes or restricted areas, is not part of the
setting studied here.
 
\subsection{Modular prediction and planning}
\label{sec:problem-pipeline}
 
In the prevailing pipeline, the decision process of the ego aircraft is
organized as two modules connected in sequence.
 
The first module is a trajectory predictor. It is a learned model
$f_{\theta}$, with trainable parameters $\theta$, that maps the observed
neighbour histories to a probabilistic forecast of where each neighbour
will be over the next $T$ steps,
\begin{equation}
	\label{eq:predictor}
	\widehat{\mathcal{T}}_t
	= f_{\theta}\big( H_t \big)
	= \big\{ q^{\,i}_{\theta}\big( \bm{p}^{i}_{t+1:t+T} \mid H^i_t \big)
	\big\}_{i \in \mathcal{N}_t},
\end{equation}
where $q^{\,i}_{\theta}$ is a predictive distribution over the future
positions of neighbour $i$. We write $\hat{\bm{p}}^{\,i}_{t+k \mid t}$
for a point summary of this distribution, read as ``the predicted
position of neighbour $i$ at time $t+k$, as forecast at time $t$,'' and
the distribution carries an uncertainty representation that may vary
with the prediction horizon and the encounter context. The concrete
parameterization of both is introduced in Section~\ref{sec:method}.
 
The second module is a planner. It is a decision rule $\pi_{\phi}$, with
design parameters $\phi$, that selects the control action from the
current state, the forecast, and the reference route,
\begin{equation}
	\label{eq:planner}
	\bu_t = \pi_{\phi}\big( \bx_t,\,
	\widehat{\mathcal{T}}_t,\,
	\bm{r}_{t:t+T} \big).
\end{equation}
The planner used in this study is \emph{separation-constrained}. Among
the actions consistent with the vehicle dynamics and actuator limits, it
uses the forecast to select actions intended to satisfy the separation
requirement~\eqref{eq:sep-requirement} over the planning horizon. Because
the requirement refers to true future positions while the planner only
holds uncertain forecasts, satisfying it involves keeping a margin around
the predicted traffic; the form of this margin used in this study is
specified in Section~\ref{sec:method}. The planner operates
in a receding-horizon manner, meaning that only the first action of each
plan is executed before the plan is refreshed with new observations at
the next step.
 
A structural feature of this pipeline matters for everything that
follows. In the modular formulation considered here, information flows
in one direction only. The predictor is trained beforehand, by comparing
its forecasts against recorded trajectories; at operation time its
output is handed to the planner, and the selected actions and realized
closed-loop outcomes do not enter the predictor's training objective.
The two
modules are also evaluated separately in current practice, the predictor
by how close its forecasts fall to recorded motion and the planner by the
behaviour of the operating loop, and the next subsection makes both
families of measures precise.
 
\subsection{Prediction accuracy and operational performance}
\label{sec:problem-objectives}
 
\paragraph{Prediction accuracy.}
The primary open-loop prediction measure considered in this study is
displacement error, the straight-line distance between where a neighbour was predicted
to be and where it actually went. Let $\mathcal{D}$ denote a collection
of prediction instances, where each instance $\zeta$ records a forecast
time $t_{\zeta}$, the observed histories at that time, and the realized
positions over the following $T$ steps; a recorded episode typically
contributes many such instances as the forecast time slides forward.
Averaging over instances, neighbours, and the prediction horizon gives
the average displacement error
\begin{equation}
	\label{eq:ade-loss}
	\mathcal{L}_{\mathrm{pred}}(\theta)
	= \frac{1}{|\mathcal{D}|} \sum_{\zeta \in \mathcal{D}}
	\frac{1}{|\mathcal{N}_{\zeta}|\, T}
	\sum_{i \in \mathcal{N}_{\zeta}} \sum_{k=1}^{T}
	\big\| \hat{\bm{p}}^{\,i}_{t_{\zeta}+k \mid t_{\zeta}}
	- \bm{p}^{i}_{t_{\zeta}+k} \big\|_2,
\end{equation}
where $\bm{p}^{i}_{t_{\zeta}+k}$ is the recorded position of neighbour
$i$ and $\mathcal{N}_{\zeta}$ is the neighbour set of instance $\zeta$;
in the single-neighbour encounters of this study,
$|\mathcal{N}_{\zeta}| = 1$.
Variants of the same idea, such as the error at the final prediction
step only, are common. These displacement-based measures are
\emph{open-loop}: they are computed by replaying recorded data, and no
controller acts on the forecast while it is being scored.
 
\paragraph{Operational performance.}
The operating system, by contrast, is judged on what happens when the
planner acts on the forecasts. Consider a mission episode $\xi$,
comprising an initial state, the motion of the neighbouring traffic, the
wind, and the reference route. Running the pipeline of
Section~\ref{sec:problem-pipeline} on this episode produces a realized
\emph{closed-loop} trajectory, so called because the aircraft's actions
feed back into what it subsequently observes. This trajectory, and every
quantity computed from it, depends jointly on the predictor parameters
$\theta$ and the planner parameters $\phi$. Three quantities summarize the
operational outcome of an episode.
 
The first is the loss-of-separation indicator
\begin{equation}
	\label{eq:phi-los}
	\Phi_{\mathrm{los}}(\theta, \phi;\, \xi)
	= \mathbb{1}\Big\{ \min_{t,\; i \in \mathcal{N}_t}
	\big\| \bm{p}_t - \bm{p}^i_t \big\|_2 < d_{\mathrm{sep}} \Big\},
\end{equation}
which equals one if the ego aircraft comes closer than the separation
standard to any neighbour at any time during the episode, and zero
otherwise. We refer to an episode containing at least one loss of
separation as a \emph{conflict episode}, and to the fraction of conflict
episodes in a set of encounters as the \emph{conflict rate}; this is the
primary safety measure used throughout the paper. Its population
counterpart is the closed-loop loss-of-separation risk
\begin{equation}
	\label{eq:conflict-risk}
	\mathcal{C}(\theta, \phi)
	= \E_{\xi \sim \mathcal{P}}\big[ \Phi_{\mathrm{los}}(\theta, \phi;\, \xi) \big],
\end{equation}
where $\mathcal{P}$ is the distribution of episodes in the target
operation.
 
The second is the terminal lateral deviation
\begin{equation}
    \label{eq:phi-dev}
    \Phi_{\mathrm{dev}}(\theta,\phi;\xi)
    =
    \frac{1}{100}
    \left\|
    P_{\perp}
    \left(
    \bm{p}_{T_{\mathrm{end}}}
    -
    \bm{r}_{T_{\mathrm{end}}}
    \right)
    \right\|_2 ,
\end{equation}
where $P_{\perp}$ projects the terminal position error onto the plane
transverse to the corridor reference axis. The factor $1/100$
normalizes the terminal lateral displacement used in Stage~2 training.
This term penalizes ending the training rollout away from the reference
route rather than the time-averaged or maximum lateral excursion over
the trajectory.

The third is the control effort
\begin{equation}
	\label{eq:phi-effort}
	\Phi_{\mathrm{effort}}(\theta, \phi;\, \xi)
	= \sum_{t=0}^{T_{\mathrm{end}}-1} \| \bu_t \|_{R}^{2},
\end{equation}
a standard quadratic measure of how hard the actuators work over the
episode, with weight matrix $R \succ 0$; it is a proxy for control
activity rather than a physical energy measurement.
 
Weighting these terms by non-negative coefficients that encode
operator preferences, the closed-loop operation of the pipeline is
summarized by the operational risk
\begin{equation}
	\label{eq:operational-risk}
	\mathcal{R}_{\mathrm{op}}(\theta, \phi)
	= \E_{\xi \sim \mathcal{P}}\!\Big[
		\lambda_{s}\, \Phi_{\mathrm{los}}
		+ \lambda_{d}\, \Phi_{\mathrm{dev}}
		+ \lambda_{u}\, \Phi_{\mathrm{effort}}
		\Big],
\end{equation}
where each term is evaluated at $(\theta, \phi;\, \xi)$. The
coefficients absorb any normalization required to place terms with
different physical units on a common scale. Further diagnostic
quantities used in the evaluation, such as the conflict-detection lead
time, are introduced together with the experimental protocol in
Section~\ref{sec:experiments}.
 
The open-loop prediction objective and the closed-loop operational
objective differ in structure, and the difference is the crux of the
paper.
$\mathcal{L}_{\mathrm{pred}}$ depends on the predictor alone and can be
computed from recorded data without ever running a controller.
$\mathcal{R}_{\mathrm{op}}$ depends on the predictor and the planner
jointly, and only through the trajectory their interaction produces. In
current practice the predictor is trained and selected against the first
measure, while the system it serves is judged by the second.
 
This mismatch raises the first question studied in the paper. For the
predictors evaluated here, and with the separation-constrained planner
held fixed at a common configuration $\bar{\phi}$, does the ordering
induced by open-loop displacement error $\mathcal{L}_{\mathrm{pred}}$
agree with the ordering induced by closed-loop loss-of-separation risk
$\mathcal{C}(\theta, \bar{\phi})$? If it does, displacement error is
a faithful proxy and the prevailing practice of selecting predictors by
that error rests on solid ground; if it does not, model selection based
on displacement error need not deliver the safest closed-loop operation.
 
The same mismatch also admits a constructive response, and this leads to
the second question. Because the closed-loop outcome depends on
$\theta$ through the planner, the training signal itself can be routed
through the planner, so that operational consequences rather than
displacement error alone shape the predictor's parameters. Whether doing
so helps, however, is a separate matter from whether the predictor has
seen the relevant encounters, since the same predictor can also be
trained on the same encounter set, under matched training conditions, by
displacement error alone. The second question is therefore whether
training the predictor through the separation-constrained planner
improves closed-loop operation beyond such matched fine-tuning that
receives no signal from the planner.
 
Answering either question requires the ability to run the same encounter
repeatedly while changing one element at a time, which recorded flight
data alone cannot provide. These two questions motivate the differentiable closed-loop framework,
the matched comparison design, and the controlled experiments developed below.
 

\section{Methodology}
\label{sec:method}
 
Answering the two questions posed at the end of
Section~\ref{sec:problem} places three requirements on the study. The
first question compares predictors by the closed-loop operation they
induce, so it requires a closed loop that can be run repeatedly on the
same encounters while the predictor is exchanged and the
separation-constrained planner is held fixed. The second question asks
whether a training signal that has passed through the planner carries
value beyond exposure to the same encounters, so it additionally
requires a differentiable path from the realized operational outcome
back to the predictor parameters, together with a comparison design
that separates the effect of the training signal from the effect of the
training data. \method is the framework built to these requirements. It
was designed under the working hypothesis that gradients routed
through the planner would steer predictive capacity toward the
situations that decide separation; the framework treats this as a hypothesis to be evaluated empirically in the closed loop.
 
\subsection{Overview}
\label{sec:method-overview}
 
\begin{figure}[!t]
	\centering
	\includegraphics[width=\textwidth]{figures/figure1_architecture_v3.pdf}
	\caption{Overview of the \method{} framework. At each control step,
	the trajectory predictor $f_{\theta}$ maps the observed neighbour
	histories to a probabilistic forecast, whose mean trajectory enters
	the planning layer. The planning layer has two levels. A convex
	quadratic program plans commanded accelerations over a short horizon
	on a double-integrator surrogate model, subject to an acceleration
	bound and to discrete-time control-barrier separation constraints
	linearized along the lines of sight to the predicted neighbour
	positions. A saturated inner control law then converts the first
	commanded acceleration into thrust and body moments. The resulting
	control acts on the full six-degree-of-freedom vehicle dynamics,
	on which wind acts as an unmeasured disturbance, and the realized
	closed-loop trajectory yields the operational quantities of
	Section~\ref{sec:problem-objectives}. In Stage~2 training, the
	task-aligned loss is back-propagated along the dashed feedback path,
	through the vehicle dynamics, the inner law, and the quadratic
	program, into the predictor parameters $\theta$; the planner
	configuration $\phi$ is fixed throughout. A separate calibration branch consumes the empirical error distribution
of that point forecast, supporting an empirical coverage check and the
Conformal-MPC comparison.}
	\label{fig:architecture}
\end{figure}
 
At each control step, \method composes three operations,
\begin{equation}
	\label{eq:method-composition}
	\widehat{\mathcal{T}}_t = f_{\theta}(H_t),
	\qquad
	\bm{a}_{t \mid t}
	= \mathrm{QP}_{\phi}\big( \bx_t,\,
	\{\hat{\bm{p}}^{\,i}_{t+k \mid t}\}_{k \leq H_p},\,
	\bm{r}_{t:t+H_p} \big),
	\qquad
	\bu_t = \mathrm{Inner}\big( \bx_t,\, \bm{a}_{t \mid t} \big),
\end{equation}
in which $f_{\theta}$ is the trajectory predictor of
Section~\ref{sec:problem-pipeline}, $\mathrm{QP}_{\phi}$ is a convex
planning program that outputs a commanded acceleration
$\bm{a}_{t \mid t} \in \R^3$ from the point summaries
$\hat{\bm{p}}^{\,i}_{t+k \mid t}$ of the forecast over a planning
horizon $H_p$, and $\mathrm{Inner}$ is a low-level control law that
converts the commanded acceleration into the thrust and body-moment
command $\bu_t$ applied to the vehicle. The planner configuration
$\phi$ collects the cost weights, the barrier coefficient, and the
acceleration bound of the program. It is held fixed across the
mainline predictor comparisons and the matched Stage 1b--Stage 2
analysis, with alternative configurations appearing only in the
designated baseline and sensitivity analyses; how it is selected is
described in Section~\ref{sec:experiments}. Subsequent observations close the loop
in the usual receding-horizon manner, and the realized trajectory
evolves under the full vehicle dynamics~\eqref{eq:dynamics}, with wind
acting as an unmeasured disturbance.
 
The predictor answers where the surrounding traffic will be; the
planning layer selects, within the vehicle's authority, a motion that
tracks the route while keeping the required clearance from the
predicted traffic; and the vehicle dynamics render the evaluation on
the motion the aircraft actually flies. Two properties of the
composition carry the study, and each is established where the first
question or the second requires it.
Section~\ref{sec:method-planner} states the sense in which the
planning layer keeps the vehicle separated, with the same constraint construction applied to every
compared predictor;
Section~\ref{sec:method-diff} states the sense in which the
composition can be differentiated back to $\theta$;
Sections~\ref{sec:method-tasl} and~\ref{sec:method-training}
construct the loss and the matched comparison design built on them.
 
\subsection{Trajectory predictor and uncertainty calibration}
\label{sec:method-pred}
 
The predictor supplies two things to the study. The planning layer
consumes a point forecast of each neighbour's future positions. A
separate calibration branch consumes the empirical error distribution
of that point forecast, supporting the coverage analysis of
Section~\ref{sec:experiments} and one comparison arm. We describe the
network, then the two outputs in turn.
 
\paragraph{Network.}
We instantiate $f_{\theta}$ as a shared per-neighbour sequence
encoder. The observed history of a neighbour, a window of $L$ past
positions, is augmented with finite-difference velocities, so that the
encoder receives explicit motion cues rather than having to recover
them; the resulting $L \times 6$ sequence is processed by a two-layer
gated recurrent unit (GRU) with hidden width $128$, and the final
hidden state passes through a two-layer perceptron to form a context
vector. Three linear heads decode this context into, for every horizon
step $k = 1, \dots, T$ and every mixture component
$m = 1, \dots, K$, a mixture weight, a mean position, and a diagonal
log-standard-deviation, the last clamped to a numerically safe range.
The same network is applied to each neighbour independently, so the
predictive distribution of Section~\ref{sec:problem-pipeline}
factorizes per neighbour as the Gaussian mixture
\begin{equation}
	\label{eq:gmm-decoder}
	q^{\,i}_{\theta}\big( \bm{p}^{i}_{t+k} \mid H^i_t \big)
	= \sum_{m=1}^{K} \alpha^{i,m}_{t+k \mid t}\,
	\mathcal{N}\!\big( \bm{\mu}^{i,m}_{t+k \mid t},\,
	\Sigma^{i,m}_{t+k \mid t} \big),
	\qquad
	\Sigma^{i,m}_{t+k \mid t} \ \text{diagonal},
\end{equation}
with all parameters output by the network. The predictor is
conditioned only on the observed motion history of the neighbouring
aircraft. This matches
the experimental design, in which neighbouring aircraft follow replayed
trajectories whose futures do not react to the simulated wind or to the
ego aircraft, so richer conditioning would introduce inputs with no
causal relation to the prediction target. The encoder is deliberately
kept simple. The predictor is one interchangeable element of a
comparison that spans a wide range of predictive quality
(Section~\ref{sec:experiments}), and a compact, transparent network
keeps the comparison legible; the contribution of this paper does not
lie in the encoder.
 
The multi-modal head reflects that low-altitude conflict geometry
admits qualitatively different continuations. A neighbouring aircraft
approaching a crossing may turn, hold, or continue, and the mixture
can represent these branches separately. The mainline controller,
however, uses the mixture mean as a common point-forecast interface;
the modal structure and the covariance outputs do not enter the
mainline planning path. The point summary consumed by the planner is
the mixture mean trajectory,
\begin{equation}
	\label{eq:mean-traj}
	\hat{\bm{p}}^{\,i}_{t+k \mid t}
	= \bar{\bm{\mu}}^{\,i}_{t+k \mid t}
	= \sum_{m=1}^{K} \alpha^{i,m}_{t+k \mid t}\,
	\bm{\mu}^{i,m}_{t+k \mid t},
\end{equation}
which realizes the point summary introduced abstractly in
Section~\ref{sec:problem-pipeline}. The same trajectory is the
quantity that the displacement error~\eqref{eq:ade-loss} scores. That
one object serves simultaneously as the planner's input and as the
target of the conventional accuracy metric is what makes the
comparison between the two orderings of
Section~\ref{sec:problem-objectives} concrete.
 
\paragraph{Split-conformal calibration of the point-forecast error.}
For the uncertainty-aware comparison arm and the coverage analysis, we
calibrate an empirical error radius around the mixture-mean forecast;
the calibration does not use the mixture covariances and does not
alter the mainline planning path~\citep{lindemann2023safe}. The
procedure is simple to state. On a calibration set of encounters
disjoint from the training data and generated independently from the
held-out trajectory pool, the frozen predictor is run, and the
non-conformity score of each prediction instance is the Euclidean
error of the mean trajectory,
$s = \| \hat{\bm{p}}^{\,i}_{t+k \mid t} - \bm{p}^{i}_{t+k} \|_2$,
collected over the look-ahead steps within the planning horizon. For a
user-specified miscoverage level $\delta \in (0,1)$, the conformal
radius is the finite-sample quantile
\begin{equation}
	\label{eq:conformal-radius}
	r_{\mathrm{conf}}(\delta)
	= \mathrm{Quantile}\!\left(
	\{ s_j \}_{j=1}^{n_{\mathrm{cal}}};\;
	\frac{\lceil (n_{\mathrm{cal}} + 1)(1 - \delta) \rceil}
	{n_{\mathrm{cal}}}
	\right),
\end{equation}
where the corrected quantile level, rather than the plain
$1 - \delta$, is the standard small-sample adjustment of split
conformal prediction. Under the exchangeability assumption between the
calibration and deployment data, the radius carries the standard
marginal coverage interpretation: with probability at least
$1 - \delta$, the true neighbour position lies within
$r_{\mathrm{conf}}$ of the predicted mean. The guarantee is marginal
over the pooled neighbour--horizon prediction instances used for
calibration; it is not a simultaneous trajectory-level coverage
statement, and because scores drawn from the same encounter are
mutually dependent, the radius is read as an empirical pooled quantile
carrying this marginal interpretation. Two features of this study make that assumption unusually clean.
The neighbour trajectories are replayed rather than reactive, so the
distribution of prediction targets does not depend on the ego policy,
and the planner configuration is fixed, so no feedback from training
alters the deployment distribution.
 
Calibration plays two roles in the study. First, it provides an
empirical check of whether the target marginal coverage level is
attained on independently generated evaluation encounters. Second, it
defines the Conformal-MPC comparison, in which the same planning
formulation is run with its separation threshold enlarged to
$d_{\mathrm{sep}}+r_{\mathrm{conf}}$, while conflicts are still judged
at the true operational threshold $d_{\mathrm{sep}}$. The mainline
planner used by the other configurations enforces separation against
the predicted mean with the fixed buffer $d_{\mathrm{sep}}$; forecast
uncertainty therefore enters through the calibration analysis and the
Conformal-MPC comparison, not through the mainline control path.
 
\subsection{Separation-constrained planning layer}
\label{sec:method-planner}
 
Section~\ref{sec:problem-pipeline} left the planner abstract, as a
rule that selects actions intended to satisfy the separation
requirement~\eqref{eq:sep-requirement} against the predicted traffic.
This subsection gives the layered realization used throughout the
study. Planning is performed on a simplified motion model by a convex
program that can be solved reliably at every control step and, for the
purposes of Section~\ref{sec:method-diff}, differentiated; the
program's output is then converted to vehicle commands by a low-level
control law; and the commands act on the full vehicle dynamics, which
appear only in the closed-loop rollout. Planning on a surrogate model
while the plant runs at full fidelity is standard practice in
hierarchical control, and it has a useful side effect here. The
surrogate exposes the vehicle's manoeuvring authority as a single
interpretable bound, which the experiments later vary as a controlled
factor.
 
\paragraph{From clearance to a barrier recursion.}
The most direct encoding of separation would require, at every
planned step, that the ego position stay outside each neighbour's
keep-clear region. Imposed step by step this is brittle. A vehicle
that is merely outside the region now can still be on a closing
course, discovering the conflict only when avoidance is no longer
within its authority. A discrete-time control barrier function
constrains not only the current clearance but the rate at which
clearance may shrink, so the vehicle is obliged to begin reacting
while a correction is still cheap. Concretely, let $h \geq 0$ denote a
clearance measure defined below. The barrier recursion
\begin{equation}
	\label{eq:cbf-condition}
	h^{i}_{k} \;\geq\; (1 - \gamma)\, h^{i}_{k-1} - \epsilon^{i}_{k},
	\qquad \gamma \in (0, 1],\quad \epsilon^{i}_{k} \geq 0,
\end{equation}
requires the clearance at step $k$ to remain at least a fraction
$(1-\gamma)$ of its previous value, so that the margin may shrink by
at most a fraction $\gamma$ per step, up to a slack $\epsilon^i_k$
whose role is discussed below. The coefficient $\gamma$ sets how quickly the
vehicle may consume its safety margin. A smaller $\gamma$ imposes a
more conservative evolution of the clearance, whereas a larger
$\gamma$ permits faster margin reduction. In our implementation
$\gamma$ is a fixed constant of the compiled program. The
differentiable solver compiles the optimization under
disciplined-parametrized-programming rules; treating $\gamma$ as a
program parameter would multiply it by expressions that already depend
on other parameters through the recursion, which those rules forbid,
so $\gamma$ is fixed before compilation. The value of $\gamma$, along
with the other planner constants, is fixed before training; the
selection protocol is described in Section~\ref{sec:experiments}.
 
\paragraph{Clearance against the predicted traffic.}
The clearance measure is built from the predicted mean positions. For
neighbour $i$ and planned step $k$, let
$\bm{n}^{i}_{k}$ be the unit vector from the predicted position
$\hat{\bm{p}}^{\,i}_{t+k \mid t}$ toward the corresponding
reference-route point, a linearization direction along which a local
supporting half-space of the keep-clear region is constructed. The clearance is the
linearized signed distance
\begin{equation}
	\label{eq:cbf-barrier}
	h^{i}_{k}
	= \bm{n}^{i\,\top}_{k}
	\big( \bm{p}_{t+k \mid t} - \hat{\bm{p}}^{\,i}_{t+k \mid t} \big)
	- d_{\mathrm{sep}},
\end{equation}
where $\bm{p}_{t+k \mid t}$ is the planned ego position. Because a
unit-direction projection never exceeds the Euclidean norm,
$h^i_k \geq 0$ implies
$\| \bm{p}_{t+k \mid t} - \hat{\bm{p}}^{\,i}_{t+k \mid t} \|_2 \geq
	d_{\mathrm{sep}}$; the constraint therefore keeps the planned
position in a supporting half-space of the keep-clear ball around the
predicted mean, a conservative linear sufficient condition for the
separation requirement against that mean. The linearization is what
renders every separation constraint affine in the decision variables,
so the planning program below remains a convex quadratic program.
 
\paragraph{The planning program.}
At each control step $t$, the upper layer solves, on a
double-integrator surrogate with planned positions $\bm{p}_k$,
velocities $\bm{v}_k$, and commanded accelerations $\bm{a}_k$ (we
abbreviate $\bm{p}_{t+k \mid t}$ as $\bm{p}_k$ inside the program, and
likewise for $\bm{v}$ and $\bm{a}$),
\begin{equation}
	\label{eq:planning-qp}
	\begin{aligned}
		\min_{\bm{a},\, \bm{p},\, \bm{v},\, \bm{\epsilon}} \quad
		 & w_{\mathrm{track}} \sum_{k=0}^{H_p}
		\| \bm{p}_k - \bm{r}_{t+k} \|_2^2
		+ w_{\mathrm{acc}} \sum_{k=0}^{H_p - 1} \| \bm{a}_k \|_2^2
		+ w_{\mathrm{slack}} \sum_{i,k} \epsilon^{i}_{k}          \\
		\text{s.t.} \quad
		 & \bm{p}_0 = \bm{p}_t, \quad \bm{v}_0 = \bm{v}_t,        \\
		 & \bm{p}_{k+1} = \bm{p}_k + \Delta t\, \bm{v}_k, \quad
		\bm{v}_{k+1} = \bm{v}_k + \Delta t\, \bm{a}_k,
		\quad k = 0, \dots, H_p - 1,                              \\
		 & \| \bm{a}_k \|_{\infty} \leq a_{\max},
		\quad k = 0, \dots, H_p - 1,                              \\
		 & h^{i}_{k} \geq (1 - \gamma)\, h^{i}_{k-1}
		- \epsilon^{i}_{k}, \qquad
		\epsilon^{i}_{k} \geq 0,
		\quad i \in \mathcal{N}_t,\; k = 1, \dots, H_p,
	\end{aligned}
\end{equation}
with $h^i_0$ evaluated at the current ego and neighbour positions,
under the convention $\hat{\bm{p}}^{\,i}_{t \mid t} := \bm{p}^i_t$
that the zeroth forecast step is the currently observed position. At
the boundaries of the horizon, $h^i_0$ enters only the first
inequality and $h^i_{H_p}$ only the last, while each interior
clearance $h^i_k$ enters the inequalities at steps $k$ and $k+1$. The program consumes the first $H_p$ steps of the $T$-step
forecast, with the final predicted position held when the planning
horizon exceeds the available forecast. The objective reads as the
operator's own trade-off. Track the assigned route, spend little
control effort, and respect the barrier recursion against every
neighbour, all within the acceleration envelope $a_{\max}$ that
summarizes the vehicle's manoeuvring authority on the surrogate model.
The program is a single convex quadratic program per control step,
solved through a differentiable convex-optimization
layer~\citep{mandi2024decisionfocused}; its solution map is written
\begin{equation}
	\label{eq:solution-map}
	\bm{a}_{t \mid t}
	= \mathrm{QP}_{\phi}\big( \bx_t,\,
	\{\hat{\bm{p}}^{\,i}_{t+k \mid t}\}_{k \leq H_p},\,
	\bm{r}_{t:t+H_p} \big),
\end{equation}
where $\phi = (w_{\mathrm{track}}, w_{\mathrm{acc}},
	w_{\mathrm{slack}}, \gamma, a_{\max}, H_p)$ collects the fixed
configuration.
 
The slack variables carry operational meaning that recurs in the
analysis. A hard barrier recursion can become infeasible, most often
when a predicted position closes in faster than the acceleration
envelope can answer, and a planner with no recourse would fail to
return a command. The slacks restore feasibility at a large linear penalty
$w_{\mathrm{slack}}$. Under standard exact-penalty conditions and a
sufficiently large weight, the softened program recovers a zero-slack
solution whenever the hard program is feasible; we do not assume that
this condition has been verified for every encounter. Positive slack is therefore read as a diagnostic of barrier pressure
under the chosen objective and envelope, rather than as proof that no
zero-slack solution exists; the later residual-conflict attribution
uses it only in this diagnostic role.
 
\paragraph{Inner control law.}
The program's first commanded acceleration is converted to vehicle
commands by a saturated feedback law. The commanded inertial
acceleration, with gravity compensated, determines a desired thrust
magnitude; the lateral components of the desired force determine
target roll and pitch angles, clamped to a physical tilt envelope; and
a proportional--derivative attitude loop converts the angle errors
into body moments, clamped to the vehicle's moment limits. The
resulting thrust-and-moment command $\bu_t$ is what acts on the full
dynamics~\eqref{eq:dynamics}. Together with the acceleration bound of the program, these limits
define the physical envelope within which avoidance must be
accomplished.
 
\paragraph{What the constraint construction guarantees.}
The following statement records the safety property of the upper
layer. It is deliberately stated for a single planning instant and on
the planning model, and its limitations are discussed immediately
after.
 
\begin{proposition}[Forward invariance of the planned clearance]
	\label{prop:forward-invariance}
	Fix a planning instant $t$ and the predicted mean positions
	$\{\hat{\bm{p}}^{\,i}_{t+k \mid t}\}$ available at that instant.
	Suppose the program~\eqref{eq:planning-qp} is feasible without
	slack ($\bm{\epsilon} = 0$) and $\gamma \in (0, 1]$. Then for every
	neighbour $i$ with initial clearance $h^i_0 \geq 0$, the planned
	trajectory satisfies $h^i_k \geq 0$ for all $k \leq H_p$, and
	consequently
	$\| \bm{p}_{t+k \mid t} - \hat{\bm{p}}^{\,i}_{t+k \mid t} \|_2
		\geq d_{\mathrm{sep}}$ at every planned step.
\end{proposition}
 
\noindent\emph{Proof sketch.}
Unrolling the recursion~\eqref{eq:cbf-condition} with
$\bm{\epsilon} = 0$ gives $h^i_k \geq (1-\gamma)^k h^i_0 \geq 0$, and
$h^i_k \geq 0$ implies the Euclidean separation by the half-space
bound noted at~\eqref{eq:cbf-barrier}. The full argument is in the
appendix.\hfill$\square$
 
When the recursion is feasible without slack, the construction thus
certifies that the planned motion keeps the required clearance from
every predicted position over the horizon. We refer to the constraint
construction in this role as a separation certificate, and to
operation in which it remains feasible as certificate-governed; both
terms are used in this conditional sense throughout the paper.
 
Three qualifications delimit what is certified, and each maps onto a
quantity the experiments measure. First, the certificate binds the
planned trajectory to the currently predicted positions; separation
from the true traffic holds insofar as the prediction is adequate at
the planned passing geometry, which is precisely the coupling between
predictive quality and safety that the study interrogates rather than
assumes. Second, the planner is executed in receding horizon, so the
barriers are rebuilt at every step around refreshed predictions; the
single-instant certificates chain into a closed-loop property only
when each replanned program inherits a feasible, non-negative
clearance from its predecessor, and when a refreshed prediction moves
the keep-clear region faster than the envelope can answer, that
inheritance may fail and positive slack may appear. Third, the certificate
lives on the double-integrator surrogate, while the aircraft flies the
full dynamics through the saturated inner law; surrogate--plant
mismatch and the tilt and moment clamps bound how faithfully the
planned motion is realized. The experiments therefore report
separation on the realized trajectory, attribute residual conflicts
through the slack signal, and treat the certificate as a per-instant
construction rather than as a global claim.
 
\subsection{Differentiating through the planning layer}
\label{sec:method-diff}
 
The second question of Section~\ref{sec:problem} requires a training
signal that originates in the realized closed-loop outcome and reaches
the predictor parameters. Two links in the chain need care, and each
rests on one condition stated here.
 
The first link is the rollout itself. The training signal is measured
on the realized trajectory, so its gradient must pass backward through
the closed-loop dynamics.
 
\begin{assumption}[Smooth vehicle dynamics]
	\label{ass:dynamics}
	The dynamics map $g$ in~\eqref{eq:dynamics} is continuously
	differentiable in the state and the control, with Jacobians bounded
	on the compact operating envelope of the study.
\end{assumption}
 
Under Assumption~\ref{ass:dynamics}, back-propagation through the
rollout is well defined. The inner control law of
Section~\ref{sec:method-planner} admits an almost-everywhere
derivative and transmits the signal from the realized motion to the
commanded acceleration; within saturated regions this derivative
vanishes, a reflection of the same finite actuation envelope that
bounds avoidance itself.
 
The second link is the planning program. The program is an
optimization rather than a closed-form map, but its solution still
varies smoothly with its inputs wherever the set of constraints active
at the optimum does not change. Intuitively, a small displacement of a
predicted neighbour position tilts the half-spaces
of~\eqref{eq:cbf-barrier} slightly, and the optimal accelerations move
by a correspondingly small, computable amount, obtained by
differentiating the optimality conditions that the solution satisfies.
 
\begin{proposition}[Existence of planning-layer gradients]
	\label{prop:gradient}
	Suppose the program~\eqref{eq:planning-qp} has a locally unique
	primal--dual solution at which the active constraint gradients are
	linearly independent, strict complementarity holds, the reduced Hessian is positive
	definite, and every predicted position lies at distance at least
	$\rho_{\min} > 0$ from its reference point, so that the
	linearization directions are well defined. Then the solution
	map~\eqref{eq:solution-map} is continuously differentiable in the
	predicted mean positions $\{\hat{\bm{p}}^{\,i}_{t+k \mid t}\}$ in
	a neighbourhood of the current inputs, and the derivative is
	obtained by solving the linearized Karush--Kuhn--Tucker system of
	the quadratic program at its
	solution~\citep{mandi2024decisionfocused,chen2021projection}.
\end{proposition}
 
\noindent\emph{Proof sketch.}
Under the stated regularity, the optimality system is a smooth system
of equations whose Jacobian in the primal--dual variables is
non-singular; the implicit function theorem yields differentiability
of the primal solution in the program parameters, which are themselves
smooth functions of the predicted positions
through~\eqref{eq:cbf-barrier}. Details are in the
appendix.\hfill$\square$
 
The result applies locally, away from active-set changes, degenerate
solutions, and the switching boundaries of the saturations in the
inner law. With these two links, an
operational loss evaluated on the realized trajectory admits a
well-defined gradient with respect to $\theta$, flowing through the
vehicle dynamics, the inner law, the quadratic program, and the
predictor's mean trajectory. This path is the formal content of the
design hypothesis stated in the section opening. Whether the signal
it carries has operational value beyond data exposure is the empirical
question, and the loss that generates the signal is constructed next.
 
\subsection{Task-aligned surrogate loss}
\label{sec:method-tasl}
 
The differentiable path determines where gradients can flow; the loss
determines what they encode. The natural training target would be the
operational risk~\eqref{eq:operational-risk} itself, but its safety
term is unusable as a training signal. The loss-of-separation
indicator~\eqref{eq:phi-los} is a zero--one outcome, so its gradient
is zero almost everywhere and undefined at the threshold, and training
on it would send no signal along the entire approach to a conflict. We
therefore train against a smooth separation-violation penalty in place
of the indicator, and retain the remaining terms of the risk, which
are already differentiable almost everywhere.
 
For an episode $\xi$ and the closed-loop trajectory it induces under
parameters $\theta$, define the per-step clearance
$\bar h_t = \min_{i \in \mathcal{N}_t}
	\| \bm{p}_t - \bm{p}^{i}_t \|_2 - d_{\mathrm{sep}}$, positive when
separated and negative when the threshold is breached. In the
one-neighbour training encounters used here the minimum reduces to a
single smooth clearance; with several neighbours it is differentiable
almost everywhere. The surrogate is the time-averaged logistic
proximity
\begin{equation}
	\label{eq:soft-los}
	\widetilde\Phi_{\mathrm{los}}(\theta; \xi)
	= \frac{1}{T_{\mathrm{end}}} \sum_{t=0}^{T_{\mathrm{end}} - 1}
	\sigma\big( -\beta\, \bar h_t \big),
\end{equation}
where $\sigma$ is the logistic function and $\beta > 0$ sets the
sharpness. The surrogate is near zero when the vehicle is comfortably
separated, rises smoothly as clearance shrinks, and saturates once the
threshold is breached. As $\beta \to \infty$ it approaches the
fraction of rollout steps at which the threshold is violated, a
time-averaged quantity distinct from the episode-level any-time
indicator~\eqref{eq:phi-los}; it is used because it supplies a dense
gradient over the entire approach to a conflict, rather than because
it is a consistent surrogate of the episode-level risk.
 
Aggregating over the training encounters, the task-aligned surrogate
loss is
\begin{equation}
	\label{eq:tasl}
	\mathcal{L}_{\mathrm{TASL}}(\theta)
	= \E_{\xi \sim \mathcal{D}_{\mathrm{cl}}}\!\Big[
		\lambda_{s}\, \widetilde\Phi_{\mathrm{los}}(\theta; \xi)
		+ \lambda_{d}\, \Phi_{\mathrm{dev}}(\theta; \xi)
		+ \lambda_{u}\, \Phi_{\mathrm{effort}}(\theta; \xi)
		\Big]
	+ \lambda_{a}\, \mathcal{L}_{\mathrm{ADE}}(\theta;\, \mathcal{D}_{\mathrm{cl}}),
\end{equation}
where the terminal lateral-deviation term $\Phi_{\mathrm{dev}}$
from~\eqref{eq:phi-dev} and the control effort
$\Phi_{\mathrm{effort}}$ from~\eqref{eq:phi-effort} are evaluated on
the realized trajectory and are smooth, or piecewise smooth, in the
parameters, and $\mathcal{L}_{\mathrm{ADE}}$ denotes the displacement
objective $\mathcal{L}_{\mathrm{pred}}$ of~\eqref{eq:ade-loss},
written with this subscript for emphasis. The displacement anchor $\lambda_a
	\mathcal{L}_{\mathrm{ADE}}$ serves one purpose. Without it, the
predictor could drift toward forecasts that happen to please the
planner while no longer describing the data; the anchor holds the
forecast to the recorded neighbour motion while the operational terms
act, and it is evaluated on the same encounter batches as the
operational terms. The non-negative weights
$(\lambda_s, \lambda_d, \lambda_u, \lambda_a)$ encode the operator's
trade-off, with the separation term dominant; the numerical values used in training are reported with the
experimental protocol below.
 
The contrast between $\mathcal{L}_{\mathrm{TASL}}$ and
$\mathcal{L}_{\mathrm{ADE}}$ is the design's point. Under the
displacement loss, every step of every predicted trajectory
contributes to the training signal in proportion to its error, so
predictive capacity is spent wherever error is large. Under the
task-aligned loss, the operational part of the signal reaches the
predictor only through the planner, and the planner is sensitive to
the forecast only where its separation constraints bind. This
localization can be stated precisely.
 
\begin{proposition}[Support of the operational gradient]
	\label{prop:gradient-support}
	Under Assumption~\ref{ass:dynamics} and the regularity conditions
	of Proposition~\ref{prop:gradient}, the gradient of the operational
	part of~\eqref{eq:tasl} with respect to a predicted mean position
	$\hat{\bm{p}}^{\,i}_{t+k \mid t}$ vanishes unless at least one
	barrier inequality that is active at the program optimum depends on
	the clearance $h^{i}_{k}$; through the
	recursion~\eqref{eq:cbf-condition}, these are the inequalities at
	steps $k$ and $k+1$ for neighbour $i$. Whenever no active
	inequality depends on $h^{i}_{k}$, the only training signal
	reaching that prediction is the anchor
	$\lambda_a \nabla_\theta \mathcal{L}_{\mathrm{ADE}}$.
\end{proposition}
 
\noindent\emph{Proof sketch.}
Differentiating the operational terms through the solution map gives a
factor $\partial \bm{a}_{t \mid t} / \partial
	\hat{\bm{p}}^{\,i}_{t+k \mid t}$ that, by the implicit-function form
of the optimality system, has identically zero rows for parameters appearing only in inactive inequality constraints. A predicted position enters the program only through the clearance $h^{i}_{k}$, which appears in the inequalities at steps $k$ and $k+1$, so the gradient vanishes whenever both are inactive. The full argument is in the appendix.\hfill$\square$
 
The direction of the surviving signal is determined by the local
geometry of the active linearized constraints, and it is worth stating
exactly. At an active pair, the predicted position enters the
clearance~\eqref{eq:cbf-barrier} both directly and through the
linearization direction, and with $\rho$ the range from the predicted
position to its reference point, the derivative of the clearance in
the predicted position is
$-\bigl( \bm{n}^{i}_{k}
	+ (I - \bm{n}^{i}_{k} \bm{n}^{i\,\top}_{k})
	(\bm{p}_{t+k \mid t} - \hat{\bm{p}}^{\,i}_{t+k \mid t}) / \rho
	\bigr)$,
the linearization direction plus a transverse correction proportional
to the orthogonal offset of the planned relative position. The extent to which the resulting parameter updates concentrate along
the ego--neighbour passing geometry is therefore an empirical question
examined later through the directional prediction-space analysis.
 
The proposition establishes where the operational signal can act, and
nothing about whether acting there helps. The design hypothesis that
motivated the construction goes one step further, and we record it as
a hypothesis rather than a result.
 
\begin{hypothesis}[Accuracy reallocation]
	\label{hyp:reallocation}
	Relative to a predictor trained on the same encounters under the
	displacement loss alone, training under
	$\mathcal{L}_{\mathrm{TASL}}$ reduces the prediction error of the
	mean trajectory more at prediction locations associated with active
	barrier inequalities than elsewhere, and this concentration of
	accuracy where separation is decided lowers the closed-loop
	operational risk.
\end{hypothesis}
 
The hypothesis is tested empirically by comparing both the spatial distribution of prediction error and the closed-loop outcomes of the matched training schemes defined next.
 
\subsection{Training schemes and matched comparison design}
\label{sec:method-training}
 
The second question of Section~\ref{sec:problem} asks whether
training the predictor through the planner adds operational value
beyond training the same predictor on the same encounters without the
planner's signal. Answering it requires more than the differentiable
path; it requires a comparison in which the training signal is the
only element that varies. We therefore define three training schemes
over a common predictor architecture and initialization,
\begin{equation}
	\label{eq:training-schemes}
	\theta^{(1)}
	\in \arg\min_{\theta}\,
	\mathcal{L}_{\mathrm{S1}}\big(\theta;\, \mathcal{D}_{\mathrm{hist}}\big),
	\qquad
	\theta^{(1\mathrm{b})}
	\in \arg\min_{\theta}\,
	\mathcal{L}_{\mathrm{ADE}}\big(\theta;\, \mathcal{D}_{\mathrm{cl}}\big),
	\qquad
	\theta^{(2)}
	\in \arg\min_{\theta}\,
	\mathcal{L}_{\mathrm{TASL}}\big(\theta;\, \mathcal{D}_{\mathrm{cl}}\big),
\end{equation}
where $\mathcal{D}_{\mathrm{hist}}$ denotes the historical trajectory
corpus and $\mathcal{D}_{\mathrm{cl}}$ the stream of closed-loop
training encounters, identical for the second and third schemes. The
Stage 1 objective $\mathcal{L}_{\mathrm{S1}}$ combines the mixture
negative log-likelihood with a best-mode displacement term, the
standard recipe for training multi-modal predictors; it supplies a
training signal to all three output heads of
Section~\ref{sec:method-pred}, including the mixture weights and the
covariances. Stage 1b starts from the Stage 1 parameters and
minimizes, on the closed-loop encounter stream, exactly the
mixture-mean displacement objective $\mathcal{L}_{\mathrm{ADE}}$ that
serves as the anchor term of the task-aligned loss.
Stage 2 starts from the same parameters and trains on the same stream
under the task-aligned loss, with the gradient routed through the
planner. The argmin notation is conceptual; in practice each scheme is
a finite training run under a matched protocol, and Stage 1b and Stage
2 are matched in initialization, optimizer state, encounter order,
and update budget, so that their intended difference is the training objective; their prediction-side terms are the same mixture-mean displacement objective, so the pair differs in the operational terms alone. The Stage 1 to Stage 1b contrast, by comparison, changes both the data distribution and the form of the prediction objective, and is read as a supporting comparison only. Because the neighbouring aircraft are
replayed, the prediction targets seen by Stage 1b and Stage 2 are identical sample-for-sample.
 
\begin{table}[t]
	\centering
	\caption{The three training schemes. All schemes share the
		architecture, the Stage 1 initialization where applicable, and
		matched training conditions; the planner configuration $\phi$ is
		fixed throughout and is never updated by any scheme.}
	\label{tab:training-schemes}
	\begin{tabular}{lllcc}
		\hline
		Scheme   & Data & Loss & Planner gradient & Updates \\
		\hline
		Stage 1  & $\mathcal{D}_{\mathrm{hist}}$ (historical)
		         & $\mathcal{L}_{\mathrm{S1}}$   & no  & $\theta$ \\
		Stage 1b & $\mathcal{D}_{\mathrm{cl}}$ (encounters)
		         & $\mathcal{L}_{\mathrm{ADE}}$  & no  & $\theta$ \\
		Stage 2  & $\mathcal{D}_{\mathrm{cl}}$ (same encounters)
		         & $\mathcal{L}_{\mathrm{TASL}}$ & yes & $\theta$ \\
		\hline
	\end{tabular}
\end{table}
 
The planner parameters $\phi$ remain fixed within each training run.
The planner enters Stage~2 as a fixed differentiable operator that
routes the task gradient to $\theta$; no planner parameter is updated
by training. In the primary matched comparison, Stage-1b and Stage~2
are evaluated under the same deployment planning configuration.
Cross-configuration evaluations are introduced only as a separate
robustness analysis. Table~\ref{tab:training-schemes}
summarizes the design. Its last column is the point. The control stack is common to all
schemes, and no scheme updates anything but $\theta$; attribution to
the training objective alone is claimed for the matched pair. The
comparison of Stage 1b with Stage 2 isolates the additional effect of
the planner-routed training signal, which is the operational form of
the second question of Section~\ref{sec:problem}. The comparison of
Stage 1 with Stage 1b quantifies the combined effect of exposure to
the encounter distribution and the accompanying change of prediction
objective, a supporting contrast rather than a clean isolation. The
first question is addressed separately, by evaluating predictors of
different open-loop accuracy under the common planner configuration
specified in Section~\ref{sec:experiments}.
 
For Stage 2, the gradient that reaches the predictor is worth writing
once, because it assembles the pieces established above. Applying the
chain rule to the operational part of~\eqref{eq:tasl},
\begin{equation}
	\label{eq:stage2-gradient}
	\frac{\partial \mathcal{L}_{\mathrm{TASL}}}{\partial \theta}
	= \sum_{t}\, \sum_{i,\,k}\;
	\underbrace{\frac{\partial \mathcal{L}_{\mathrm{TASL}}}
		{\partial \bu_t}}_{\substack{\text{realized trajectory,} \\
			\text{through the rollout}}}
	\;
	\underbrace{\frac{\partial \bu_t}{\partial \bm{a}_{t \mid t}}}
	_{\text{inner law}}
	\;
	\underbrace{\frac{\partial \bm{a}_{t \mid t}}
		{\partial \hat{\bm{p}}^{\,i}_{t+k \mid t}}}
	_{\substack{\text{planner sensitivity} \\
			(\text{Prop.~\ref{prop:gradient}})}}
	\;
	\underbrace{\frac{\partial \hat{\bm{p}}^{\,i}_{t+k \mid t}}
		{\partial \theta}}_{\text{predictor Jacobian}}
	\;+\; \lambda_a
	\frac{\partial \mathcal{L}_{\mathrm{ADE}}}{\partial \theta},
\end{equation}
where the first factor accumulates through the closed-loop rollout
under Assumption~\ref{ass:dynamics}, the planner sensitivity exists by
Proposition~\ref{prop:gradient} and is non-zero only at predictions
entering at least one active barrier inequality
(Proposition~\ref{prop:gradient-support}). The equation shows in one
line how an operational signal measured on the realized trajectory is
routed back into the forecaster, and it delimits what the signal can
express. Algorithm~\ref{alg:plangrad} summarizes the procedure.
 
\begin{algorithm}[t]
	\caption{Training schemes of \method}
	\label{alg:plangrad}
	\begin{algorithmic}[1]
		\Require predictor $f_{\theta}$; fixed planner configuration
		$\phi$; historical corpus $\mathcal{D}_{\mathrm{hist}}$;
		encounter stream $\mathcal{D}_{\mathrm{cl}}$; loss weights
		$(\lambda_s, \lambda_d, \lambda_u, \lambda_a)$.
		\State \textbf{Stage 1 (offline pre-training).}
		\For{mini-batch $\mathcal{B} \subset \mathcal{D}_{\mathrm{hist}}$}
		\State $\theta \leftarrow \theta
			- \eta_1 \nabla_{\theta}
			\mathcal{L}_{\mathrm{S1}}(\theta; \mathcal{B})$
		\EndFor
		\State Save $\theta^{(1)} \leftarrow \theta$; initialize
		$\theta^{(1\mathrm{b})} \leftarrow \theta^{(1)}$ and
		$\theta^{(2)} \leftarrow \theta^{(1)}$, with the optimizer
		states of both branches initialized identically.
		\State \textbf{Stage 1b (matched fine-tuning).}
		\For{mini-batch of encounters $\mathcal{B} \subset
				\mathcal{D}_{\mathrm{cl}}$}
		\State $\theta^{(1\mathrm{b})} \leftarrow \theta^{(1\mathrm{b})}
			- \eta_2 \nabla_{\theta^{(1\mathrm{b})}}
			\mathcal{L}_{\mathrm{ADE}}(\theta^{(1\mathrm{b})}; \mathcal{B})$
		\EndFor
		\State \textbf{Stage 2 (planner-guided training).}
		\For{mini-batch of encounters $\mathcal{B} \subset
				\mathcal{D}_{\mathrm{cl}}$ (same stream as Stage 1b)}
		\For{episode $\xi \in \mathcal{B}$}
		\State Roll out the closed loop: at each step, forecast
		with $f_{\theta^{(2)}}$, solve~\eqref{eq:planning-qp}, apply
		the inner law, propagate the vehicle
		dynamics~\eqref{eq:dynamics}.
		\EndFor
		\State Compute $\mathcal{L}_{\mathrm{TASL}}(\theta^{(2)};
			\mathcal{B})$ from~\eqref{eq:tasl} on the realized
		trajectories.
		\State $\theta^{(2)} \leftarrow \theta^{(2)}
			- \eta_2 \nabla_{\theta^{(2)}}
			\mathcal{L}_{\mathrm{TASL}}(\theta^{(2)}; \mathcal{B})$
		\EndFor
		\State \Return $\theta^{(1)}, \theta^{(1\mathrm{b})},
			\theta^{(2)}$
	\end{algorithmic}
\end{algorithm}
 
Stage~2 training back-propagates through the entire rollout. The
gradient traverses the aircraft dynamics at every step, the inner
control law, and the quadratic program's implicit derivative before
reaching the predictor. The optimizer updates only the predictor
parameters; the planner parameters themselves are not learned.
Nevertheless, the fixed planner configuration enters the Stage~2
update through the differentiable solution map and therefore shapes
the gradient routed to the predictor. The experimental protocol below
specifies the encounter generation, the planning configurations used
during training and evaluation, the loss weights, and the
cross-configuration robustness checks.
 
\section{Closed-Loop Experiments and Results}
\label{sec:experiments}

\subsection{Experimental design}
\label{sec:exp-design}

The experiments run the closed loop repeatedly on controlled conflict
encounters while exchanging one element at a time. Within each controlled comparison, the relevant arms share the
simulator, the encounter stream, the evaluation harness, and every
component other than the factor under study; the four predictor arms
of the ranking comparison additionally share the same deployment
planner configuration. The section first describes the environment, the data, the
arms, and the protocol, and then reports the findings in the order the
two questions of Section~\ref{sec:problem} pose them, followed by the
analyses that establish the mechanism behind the answers and the
conditions at which the answers change.

\paragraph{Data and encounter construction.}
The experiments use the open-source Generic Urban Air Mobility
(GUAM) challenge-problem release, which provides $3{,}000$
dynamically feasible trajectories for the NASA Lift+Cruise
configuration together with a documented aircraft
model~\citep{acheson2024guam,campbell2024pyguam}. GUAM contains
single-aircraft mission trajectories rather than observed
multi-aircraft corridor interactions. We therefore use the
trajectories to construct controlled conflict encounters and do not
interpret the resulting traffic patterns as representative of
observed low-altitude operations.

The $3{,}000$ trajectories are divided into a development pool and a
held-out source pool. Trajectories $0$--$2{,}499$ are used for
predictor training and for generating the closed-loop training
encounters, while trajectories $2{,}500$--$2{,}999$ form the source
pool for calibration and evaluation. Calibration and evaluation
episodes are generated independently from this pool using separate
random seeds and independently sampled transformation parameters, so
that no episode is reused across the two sets.

Each encounter pairs the ego aircraft with one non-reactive
neighbouring aircraft following a replayed GUAM trajectory. The
trajectory is transformed so that the nominal closest approach falls
below the prescribed separation standard,
$d_{\mathrm{sep}}=30$\,m, creating a controlled geometry in which an
avoidance response is required. The neighbouring aircraft follows its
replayed trajectory throughout the encounter and does not react to
the ego aircraft. The encounter-level experiments consider one
neighbour at a time so that the interaction between prediction and
planning can be examined without additional multi-aircraft
confounding. Trajectories are resampled at $\Delta t=0.2$\,s. The
predictor receives an observation window of $L=25$ steps (5\,s) and
produces a $T=30$ step (6\,s) forecast. For the primary evaluation, each closed-loop episode spans
$20$ simulation steps ($4.0$\,s), with the nominal closest approach
placed at step~$10$ ($2.0$\,s after the episode begins). This nominal
CPA is a construction parameter of the unaltered encounter geometry;
the realized closest approach may occur at a different step once the
ego aircraft manoeuvres and responds to disturbances.

\paragraph{Closed-loop simulation.}
Recorded trajectories reveal the outcome associated with one realized
sequence of actions, but they do not provide the counterfactual
outcomes required to compare different predictors and planners on the
same encounter. Stage~2 training also requires gradients to propagate
through the resulting closed-loop trajectory. We therefore implement
the corridor environment as a differentiable simulator in PyTorch.

The ego aircraft is propagated using a six-degree-of-freedom
Lift+Cruise model parameterized from the GUAM configuration and
integrated with a fourth-order Runge--Kutta scheme at
$\Delta t=0.2$\,s. Wind is represented by a mesoscale mean flow
perturbed by a corridor-scale Gaussian random field and is not
observed by the decision system. The gust field has standard
deviation $\sigma_g=3$\,m\,s$^{-1}$, with nominal scaling
$\eta_w=0.3$, corresponding to a disturbance strength of
$0.9$\,m\,s$^{-1}$. Larger disturbance levels are examined in the
robustness analyses. The neighbouring trajectories remain replayed
and therefore do not respond to either the simulated wind or the ego
aircraft.

\begin{table}[tb]
    \centering
    \small
    \caption{Experimental configurations.}
    \label{tab:arms}
    \begin{tabular}{lllll}
        \toprule
        Configuration & Predictor & Planning layer & Data used & Purpose \\
        \midrule

        \textsc{Constant-Velocity}
        & analytic
        & CBF--MPC
        & ---
        & training-free predictor reference \\

        \textsc{Fixed-Predictor}
        & Stage 1
        & CBF--MPC
        & $\mathcal{D}_{\mathrm{hist}}$
        & conventional predict-then-plan \\

        \textsc{Stage-1b}
        & Stage 1b
        & CBF--MPC
        & $\mathcal{D}_{\mathrm{cl}}$
        & displacement-fine-tuning baseline \\

        Stage~2 (\method)
        & Stage 2
        & CBF--MPC
        & $\mathcal{D}_{\mathrm{cl}}$
        & task-aligned training \\

        \textsc{Conformal-MPC}
        & Stage 1
        & CBF--MPC, calibrated margin
        & $\mathcal{D}_{\mathrm{hist}}+\mathcal{D}_{\mathrm{cal}}$
        & uncertainty-aware planner baseline \\

        \textsc{Vanilla-MPC}
        & Stage 2
        & soft penalty, no CBF
        & $\mathcal{D}_{\mathrm{cl}}$
        & certificate ablation \\

        \textsc{Soft-IPP}
        & joint (soft)
        & soft penalty, no CBF
        & $\mathcal{D}_{\mathrm{cl}}$
        & joint training without certificate \\

        \textsc{Oracle}
        & true future
        & CBF--MPC
        & ---
        & zero-error forecast reference \\

        \bottomrule
    \end{tabular}
\end{table}

\paragraph{Experimental configurations.}
Table~\ref{tab:arms} summarizes the configurations used in the
controlled comparisons. \textsc{Constant-Velocity},
\textsc{Fixed-Predictor}, \textsc{Stage-1b}, and Stage~2 (\method)
use the same certificate-equipped deployment planner while varying
the predictor supplied to it. These four configurations form the
primary comparison used to examine whether ordering predictors by
open-loop accuracy is preserved in closed-loop conflict outcomes.

\textsc{Conformal-MPC} uses the Stage~1 predictor with the same
planning formulation but enlarges the separation threshold by a
split-conformal error radius. At miscoverage level $\delta=0.1$, the
calibration set gives $r_{\mathrm{conf}}=18.93$\,m, so the planner
uses an enlarged threshold of
$d_{\mathrm{sep}}+r_{\mathrm{conf}}$ while conflicts continue to be
evaluated at the operational $30$\,m standard. The calibration and
evaluation sets each contain $3{,}000$ pooled prediction instances.
On the independently generated evaluation encounters, the calibrated
radius attains $90.8\%$ empirical pooled coverage against the
$90\%$ target, with an evaluation-side $90$th-percentile error of
$18.65$\,m.

\textsc{Vanilla-MPC} retains the Stage~2 predictor but replaces the
CBF constraints with a hinge-type soft separation penalty using the
same planning horizon and actuation limits. \textsc{Soft-IPP} trains
the predictor through this soft planner and provides a joint-training
comparison without the explicit separation certificate.
\textsc{Oracle} supplies the planner with the true future trajectory
of the neighbouring aircraft and is used to distinguish prediction
limitations from conflicts that persist even under perfect future
information.

\paragraph{Training protocol.}
Stage~1 trains the probabilistic predictor on trajectories
$0$--$2{,}499$ using the mixture negative log-likelihood together
with a best-mode displacement term. Training uses Adam with learning
rate $8\times10^{-4}$ for $30$ epochs and yaw augmentation across
encounter orientations.

Both \textsc{Stage-1b} and Stage~2 start from the same Stage~1
checkpoint and use the same closed-loop encounter stream, encounter
order, optimizer initialization, random seed, and update budget.
\textsc{Stage-1b} performs displacement-only fine-tuning using the
mixture-mean ADE objective. Stage~2 instead minimizes the task-aligned
objective, which combines closed-loop separation, terminal lateral
deviation, and control-effort terms with a displacement anchor. Both branches
run for $70$ closed-loop iterations with learning rate
$1\times10^{-4}$, batch size $16$, and rollout length $20$ steps.
The task-aligned weights are
$(\lambda_s,\lambda_d,\lambda_u,\lambda_a)
=(3.5,0.05,0.01,0.12)$, with surrogate sharpness $\beta=0.3$.
The comparison therefore controls for initialization, encounter
exposure, optimizer setup, and training budget while contrasting
displacement-only fine-tuning with the task-aligned training
objective.

We refer to
\[
(\gamma,H_p,a_{\max})=(0.4,8,10\,\mathrm{m\,s^{-2}})
\]
as the \emph{training-time planning configuration}. It is used for
the primary Stage~2 training run and for the three independently
trained Stage~2 robustness checkpoints, and was chosen to keep the
differentiated planning problem short and numerically tractable.
Closed-loop deployment evaluation uses
\[
(\gamma,H_p,a_{\max})=(0.1,15,20\,\mathrm{m\,s^{-2}}),
\]
hereafter the \emph{deployment planning configuration}. A separate
Stage~2 checkpoint is trained under this deployment configuration.
All five Stage~2 checkpoints are evaluated under the deployment
planner, and Section~\ref{sec:results-boundary} additionally
cross-evaluates them under the training-time planner to examine how
the downstream planning configuration modulates the operational
expression of predictor differences. \textsc{Soft-IPP} starts
from the same Stage~1 initialization and uses the same encounter
stream, learning rate, batch size, seed, and update budget as the
other closed-loop fine-tuning schemes, but trains through the soft
planner.

\paragraph{Deployment planner and evaluation protocol.}
The deployment planner configuration was selected from six candidate
settings using the \method{} predictor and was subsequently held
fixed across the certificate-equipped configurations, apart from the
additional calibrated margin used by \textsc{Conformal-MPC}. The
same configuration was recovered when the candidate settings were
re-evaluated on an independently generated encounter stream
(seed $999$). Because the primary evaluation stream contributed to
the initial planner selection, the absolute conflict rates reported
below are interpreted as controlled stress-test outcomes rather than
unbiased estimates of population-level operational risk.

The primary comparison uses $n=200$ evaluation encounters with seed
$12345$, which controls the encounter sampler, wind field, and the
Python, NumPy, and PyTorch random generators; conformal calibration
uses seed $777$. Four quantities characterize each configuration.
ADE is the open-loop displacement error of the predicted mean
trajectory. Conflict rate (CR) is the percentage of encounters in
which the realized separation falls below $d_{\mathrm{sep}}$ at
least once. Mean MinSep is the average, across encounters, of the
minimum realized separation attained within each encounter. Effort is
the normalized quadratic control-activity measure defined for the
closed-loop trajectory. The terminal lateral-deviation term remains part of the task-aligned
training objective but is not used to rank configurations in the main
comparison. Maximum lateral offset is reported separately as a
supplementary trajectory-deviation diagnostic in the matched training
analysis.

The conflict outcomes are paired because every configuration is
evaluated on the same encounters. For the four predictor
configurations sharing the deployment planner, overall equality of
the paired binary outcomes is assessed using Cochran's $Q$ test,
followed by pairwise exact McNemar tests with Holm correction.
Configuration-specific conflict rates are accompanied by Wilson
$95\%$ confidence intervals. The corridor penetration analysis is
reported separately as means and standard deviations over six
independent replications.

\subsection{Open-loop prediction accuracy and closed-loop conflict outcomes}
\label{sec:results-ranking}

\begin{table}[tb]
	\centering
	\small
	\caption{Open- and closed-loop performance on the evaluation conflict encounters.}
	\label{tab:main-comparison}
	\begin{tabular}{llcccc}
		\toprule
		Configuration & Predictor
		& CR (\%) $\downarrow$
		& Mean MinSep (m) $\uparrow$
		& ADE (m) $\downarrow$
		& Effort \\
		\midrule
		\multicolumn{6}{l}{\emph{Common-planner predictor comparison}} \\
		\textsc{Constant-Velocity}
		& analytic
		& 12.0 & 47.8 & 0.83 & 52.9 \\
		\textsc{Fixed-Predictor}
		& Stage 1
		& 12.5 & 46.3 & 20.90 & 51.5 \\
		\textsc{Stage-1b}
		& Stage 1b
		& 11.5 & 47.8 & 1.84 & 52.9 \\
		\method
		& Stage 2
		& 11.0 & 47.7 & 4.32 & 52.3 \\
		\midrule
		\multicolumn{6}{l}{\emph{Calibrated-margin planner baseline}} \\
		\textsc{Conformal-MPC}
		& Stage 1
		& 11.5 & 48.8 & 20.90 & 58.1 \\
		\midrule
		\multicolumn{6}{l}{\emph{Certificate-free configurations}} \\
		\textsc{Vanilla-MPC}
		& Stage 2
		& 41.0 & 29.8 & 4.32 & 19.1 \\
		\textsc{Soft-IPP}
		& joint (soft)
		& 53.0 & 29.1 & 6.45 & 17.0 \\
		\bottomrule
	\end{tabular}
\end{table}

The primary comparison evaluates whether the ordering induced by
open-loop prediction accuracy is retained after the predictor is
placed in the closed loop. Four configurations permit this comparison:
\textsc{Constant-Velocity}, \textsc{Fixed-Predictor},
\textsc{Stage-1b}, and \method{} all use the same certificate-equipped
deployment planner and differ in the predictor supplied to it.
\textsc{Conformal-MPC} uses the same Stage~1 predictor as
\textsc{Fixed-Predictor} but modifies the planning threshold through
the calibrated margin, and is therefore treated as a planner baseline
rather than as part of the predictor ranking comparison. The two
certificate-free configurations modify the planning layer itself and
are considered in the subsequent mechanism analysis.

Across the four common-planner configurations, ADE spans a
twenty-five-fold range, from $0.83$\,m for
\textsc{Constant-Velocity} to $20.90$\,m for
\textsc{Fixed-Predictor}. The corresponding conflict rates remain
within $11.0$--$12.5\%$, a range of only $1.5$ percentage points,
while their mean minimum separations range from $46.3$ to
$47.8$\,m. These mean values do not describe the lower tail of the
separation distribution, which is examined below. Because the
encounters were deliberately constructed to require conflict
resolution, the reported rates are stress-test outcomes rather than
estimates of naturally occurring operational risk.

The paired encounter-level outcomes show the same pattern. The four
common-planner configurations produce the same binary conflict
outcome in $196$ of the $200$ encounters: $175$ remain conflict-free
under all four predictors, while $21$ result in a conflict under all
four. Only four encounters exhibit any disagreement across the
predictor configurations. This high episode-level concordance is
consistent with a shared-difficulty structure among the encounters;
the source of the residual conflicts is examined in the subsequent
mechanism analysis.

Formal paired tests provide no evidence of an overall difference in
conflict probability among the four configurations. Cochran's $Q$
test gives $Q=5.00$ ($\mathrm{df}=3$, $p=0.17$), and none of the six
pairwise exact McNemar tests is significant after Holm correction.
Relative to \method{}, the discordant conflict counts are $2/0$ for
\textsc{Constant-Velocity} ($p=0.50$), $3/0$ for
\textsc{Fixed-Predictor} ($p=0.25$), and $2/1$ for
\textsc{Stage-1b} ($p=1.00$). The Wilson $95\%$ intervals for the
configuration-specific conflict rates are $[7.4,16.1]\%$ for
\method{}, $[7.8,16.7]\%$ for \textsc{Stage-1b},
$[8.2,17.2]\%$ for \textsc{Constant-Velocity}, and
$[8.6,17.8]\%$ for \textsc{Fixed-Predictor}. These results should not
be interpreted as evidence of equivalence. Rather, within the
resolution of the present experiment, they do not support
preservation of the ADE-based predictor ranking in closed-loop
conflict outcomes under the common planner.

\begin{figure}[tb]
	\centering
	\includegraphics[width=0.8\linewidth]
	{figures/fig01_ade_cr_decoupling.pdf}
	\caption{Open-loop prediction error and closed-loop conflict rate.}
	\label{fig:decoupling-scatter}
\end{figure}

Figure~\ref{fig:decoupling-scatter} makes the contrast between
open-loop prediction error and closed-loop conflict outcomes explicit.
\textsc{Constant-Velocity} attains the lowest ADE in the evaluated
encounters but shows no detectable conflict-rate advantage under the
common planner. Its low ADE is specific to the motion characteristics
of the present encounter set and should not be interpreted as a
general advantage of constant-velocity prediction.

The pattern is not driven solely by the high-error
\textsc{Fixed-Predictor}. Excluding that configuration still leaves an
approximately five-fold ADE range ($0.83$--$4.32$\,m) among the
remaining common-planner configurations, whose conflict rates differ
by no more than one percentage point. In particular,
displacement-based fine-tuning reduces the ADE of \textsc{Stage-1b}
to $1.84$\,m, compared with $4.32$\,m for \method{}, without a
detectable difference in conflict rate. Thus, within the $0.83$--$20.90$\,m ADE range spanned by the evaluated
common-planner predictors, improvements in geometric prediction
accuracy do not translate into a measurable reduction in
threshold-crossing probability under the common planner.
The same-predictor contrast that changes the planning layer rather
than the predictor is examined next.

\begin{figure}[tb]
	\centering
	\includegraphics[width=0.8\linewidth]
	{figures/fig02_minsep_ecdf.pdf}
	\caption{Empirical cumulative distributions of realized minimum separation.}
	\label{fig:minsep-cdf}
\end{figure}

Figure~\ref{fig:minsep-cdf} expands the binary conflict indicator into
the full distribution of realized minimum separation. Because a
conflict is defined by $\mathrm{MinSep}<30$\,m, the empirical
cumulative probability immediately below the $30$\,m standard
corresponds to the conflict rate. Around this threshold, the lower
tails of the CBF-constrained configurations are closely aligned,
consistent with their similar conflict rates. Descriptive differences
remain above the threshold, most visibly in the somewhat left-shifted
central portion of the \textsc{Fixed-Predictor} distribution,
consistent with its lower mean minimum separation in
Table~\ref{tab:main-comparison}. In the current zoomed view, the two
certificate-free threshold crossings lie above the displayed
ordinate range.

The certificate-free configurations place substantially more
probability below the $30$\,m standard, whereas some residual
violations under the CBF-constrained configurations extend to minimum
separations of approximately $11$\,m. The former shows that the
closed loop remains sensitive to changes outside the predictor, while
the latter shows that the common constraint layer does not eliminate
all difficult encounters. These observations motivate the analysis
of the planning layer and the attribution of residual conflicts in
the following subsection.

\subsection{Planning constraints, manoeuvring authority, and residual-conflict attribution}
\label{sec:results-mechanism}

The common-planner comparison showed that large differences in
prediction accuracy do not produce detectable differences in conflict
rate under the deployment planner. This leaves two related questions:
whether the planning layer itself accounts for the much larger
variation in closed-loop outcomes, and why a non-negligible set of
conflicts remains even when the separation constraint is active. We
address these questions by first changing the planning regime while
holding the predictor fixed, then examining the feasibility structure
of the residual conflicts and their response to perfect future
information. A planner-parameter sweep subsequently tests whether the
same residual behaviour is sensitive to the available manoeuvring
authority, and plant--model mismatch experiments assess whether these
patterns persist when the realized aircraft dynamics depart from the
planner model.

\paragraph{Effect of the separation constraint.}

To separate the effect of the planning layer from that of the
predictor, we first compare \method{} with \textsc{Vanilla-MPC}.
Both configurations use the same Stage~2 predictor and therefore have
the same ADE of $4.32$\,m; the principal change is that
\textsc{Vanilla-MPC} replaces the CBF separation constraints with a
soft repulsion penalty. Removing the CBF constraints increases the
conflict rate from $11.0\%$ to $41.0\%$, while mean minimum
separation decreases from $47.7$ to $29.8$\,m. This contrast is much
larger than the variation observed when the predictor is exchanged
behind the common CBF planner.

The joint-training configuration \textsc{Soft-IPP}, which also uses
the soft planning regime, remains in the same high-conflict region,
with a conflict rate of $53.0\%$. Because \textsc{Soft-IPP} differs
from \textsc{Vanilla-MPC} in both predictor training and the
closed-loop training objective, it is not interpreted as a matched
ablation of joint training. Rather, it provides an additional
certificate-free reference showing that the high conflict rate is not
specific to a single soft-planner configuration.

\begin{figure}[tb]
	\centering
	\includegraphics[width=0.8\linewidth]{figures/fig03_effort_cr.pdf}
	\caption{Control effort and conflict rate across planning configurations.}
	\label{fig:effort-cr}
\end{figure}

Figure~\ref{fig:effort-cr} shows the corresponding relationship
between control activity and conflict rate. The CBF-constrained
configurations cluster at conflict rates of $11.0$--$12.5\%$ with
mean control effort between approximately $51$ and $58$, whereas
\textsc{Vanilla-MPC} and \textsc{Soft-IPP} have substantially lower
effort, $19.1$ and $17.0$, together with conflict rates of $41.0\%$
and $53.0\%$. The lower effort of the certificate-free configurations
should therefore not be interpreted as an efficiency gain. It is
associated with substantially higher conflict rates and is consistent
with less active avoidance. \textsc{Conformal-MPC}, by contrast,
incurs the highest mean effort among the CBF-constrained
configurations because of its enlarged planning margin, without a
detectable reduction in conflict rate relative to the common-planner
predictor configurations.

The same-predictor ablation therefore identifies the planning regime
as a major source of the observed conflict-rate variation. It does
not, however, explain the residual conflicts that remain under the
CBF-constrained planner. We examine those episodes next.

\paragraph{Attribution of residual conflicts.}

The \method{} configuration contains $22$ conflict episodes among the
$200$ evaluation encounters. A remaining possibility is that these
episodes arise because prediction errors drive the closed loop into
states from which the planner can no longer maintain separation. We
use two complementary diagnostics to distinguish this explanation
from limitations imposed by the planning and manoeuvring envelope.

The first diagnostic evaluates the feasibility pressure encountered
during the realized Stage~2 rollouts. For each of the $22$ conflict
episodes and each of its $20$ rollout steps, the planning problem is
re-solved while holding the realized prediction stream fixed and
minimizing the barrier relaxation required to restore feasibility. A
step is classified as hard-infeasible when the minimum required
relaxation exceeds $1$\,mm.

\begin{figure}[p]
	\centering
	\begin{subfigure}{\linewidth}
		\centering
		\includegraphics[width=\linewidth]{figures/fig05b_raster.pdf}
		\caption{Minimum required barrier relaxation.}
		\label{fig:attr-relaxation}
	\end{subfigure}\\[6pt]
	\begin{subfigure}{\linewidth}
		\centering
		\includegraphics[width=\linewidth]{figures/fig05c_oracle_depth.pdf}
		\caption{Violation depth under Stage~2 and oracle replay.}
		\label{fig:attr-oracle}
	\end{subfigure}
	\caption{Feasibility and oracle analysis of residual conflicts.}
	\label{fig:attribution}
\end{figure}

Of the $440$ audited rollout steps, $414$ require more than
$1$\,mm of barrier relaxation and are therefore classified as
hard-infeasible. The result is insensitive to the precise numerical
threshold over a wide range: the count remains $414$ for thresholds
from $10^{-6}$ to $10^{-2}$\,m because no required relaxation falls
within that interval. The violations are also substantial in
magnitude. Of the $414$ hard-infeasible steps, $406$ require more
than $1$\,m of relaxation and $139$ require more than $10$\,m.

Within every residual-conflict episode, the hard-infeasible steps form
a single contiguous interval beginning at the first rollout step and
covering, on average, $18.8$ of the $20$ steps. The closest-approach
step lies inside this interval in all $22$ episodes. Feasibility is
recovered only near the end of those episodes whose infeasible
intervals are shorter than the full rollout. These results show that
the residual conflicts are not associated with isolated or marginal
constraint violations; under the realized prediction stream, the
planner operates for most of the encounter in a region where the hard
barrier conditions cannot be satisfied within the specified planning
and manoeuvring envelope.

This feasibility audit is conditional on the realized prediction
stream. Because earlier prediction errors may have influenced the
states from which later planning problems are solved, the audit alone
cannot determine whether a more accurate predictor could have avoided
the infeasible region. The second diagnostic therefore replays the
same $22$ episodes from their initial conditions while replacing the
Stage~2 forecast with the true future trajectory of the neighbouring
aircraft.

None of the $22$ residual conflicts is resolved by the oracle replay.
The oracle rollout itself requires active barrier relaxation at the
closest-approach step in $20$ of the $22$ episodes; the remaining two
have required relaxations below $10^{-6}$\,m at that step but still
experience a loss of separation elsewhere in the rollout. Figure
\ref{fig:attr-oracle} gives the corresponding violation depths,
defined as
$d_{\mathrm{sep}}-\mathrm{MinSep}$. The mean absolute change in
violation depth between the Stage~2 and oracle replays is only
$0.63$\,m, and no episode reaches the zero-violation rescue line.
Across the full $200$-encounter set, the oracle reference yields
$24$ conflicts, compared with $22$ under Stage~2. None of the
$22$ Stage~2 residual conflicts is resolved by the oracle replay;
instead, two different encounters conflict only under the oracle
reference. The oracle is therefore used here as a zero-forecast-error
diagnostic rather than as a guaranteed closed-loop safety upper bound.
In this evaluation, replacing the learned forecast with the true future
trajectory does not monotonically improve the closed-loop conflict
outcome.

Taken together, the feasibility and oracle analyses shift the
explanation of the residual conflicts away from prediction error
alone and toward the limits of the available planning and manoeuvring
envelope. The next analysis tests this interpretation directly by
varying the available acceleration authority.

\paragraph{Sensitivity to manoeuvring authority.}

If the residual conflicts are constrained by the manoeuvring envelope,
their frequency should respond strongly to the amount of corrective
acceleration available to the planner. We therefore evaluate the
Stage~2 predictor over a $4\times4$ grid of barrier coefficients
$\gamma$ and acceleration limits $a_{\max}$, using $n=200$
encounters per grid point and holding the planning horizon at
$H_p=15$.

\begin{figure}[tb]
	\centering
	\includegraphics[width=0.98\linewidth]{figures/fig11_landscape.pdf}
	\caption{Conflict-rate sensitivity to barrier and acceleration-authority settings.}
	\label{fig:planner-heatmap}
\end{figure}

The resulting surface is far more sensitive to $a_{\max}$ than to
$\gamma$. Increasing the available acceleration authority from
$5$ to $20$\,m\,s$^{-2}$ reduces the conflict rate by
$59.0$--$70.5$ percentage points across the four values of
$\gamma$, with a $70.5$-point reduction at the deployment value
$\gamma=0.1$. By comparison, the span across $\gamma$ at fixed
$a_{\max}$ is at most $10.5$ percentage points when
$a_{\max}=5$\,m\,s$^{-2}$ and at most $2.5$ percentage points at
the higher authority levels.

The deployment configuration
$(\gamma,a_{\max})=(0.1,20\,\mathrm{m\,s^{-2}})$ is one point on
this $H_p=15$ surface. Re-evaluating that point on an independently
generated encounter stream gives a conflict rate of $12.0\%$,
compared with $11.0\%$ on the primary evaluation stream, a
difference of one percentage point, or two encounters. The
training-time configuration uses $H_p=8$ and therefore does not lie
on the same surface even though its $(\gamma,a_{\max})$ coordinates
can be shown for reference.

The parameter sweep provides a direct complement to the oracle
analysis. The CBF constraint substantially changes the outcome for a
fixed predictor and acceleration envelope, while the acceleration
sweep shows that the achievable closed-loop outcome is itself strongly
constrained by the available manoeuvring authority. We therefore use
the term \emph{manoeuvring-authority-limited} for the residual
conflicts in the restricted operational sense that they remain
unresolved under oracle prediction within the fixed planning and
control envelope evaluated here. The term does not imply that no
collision-free trajectory could exist under a different planning
horizon, control architecture, or aircraft capability.

\paragraph{Robustness to plant--model mismatch.}

The preceding analyses use a planner whose surrogate dynamics and
available authority are calibrated to the nominal aircraft model. In
operation, however, mass, thrust effectiveness, response delay, and
environmental disturbances may depart from their nominal values. We
therefore perturb the simulated aircraft while keeping the planner
model and its parameters unchanged. This test asks whether the
advantage of the CBF-constrained planning regime persists when the
commands optimized by the planner are executed by a mismatched plant.

\begin{table}[tb]
	\centering
	\small
	\caption{Conflict rates under plant--model mismatch.}
	\label{tab:mismatch}
	\begin{tabular}{lccc}
		\toprule
		Regime
		& \method{} (CBF)
		& \textsc{Fixed-Predictor} (CBF)
		& \textsc{Vanilla-MPC} (no CBF) \\
		\midrule
		nominal
		& 11.0 & 12.5 & 41.0 \\
		mass $+20\%$
		& 17.5 & 18.5 & 75.0 \\
		mass $-15\%$
		& 9.0 & 10.0 & 19.5 \\
		inertia $+30\%$
		& 12.0 & 12.0 & 42.5 \\
		thrust efficiency $0.85$
		& 16.5 & 18.5 & 68.5 \\
		actuator delay $1$ step
		& 18.5 & 22.0 & 54.0 \\
		actuator delay $2$ steps
		& 28.0 & 34.0 & 87.5 \\
		wind shift
		& 11.0 & 12.0 & 39.5 \\
		combined four-factor shift
		& 30.0 & 40.0 & 94.5 \\
		\bottomrule
	\end{tabular}
\end{table}

The wind-shift condition in Table~\ref{tab:mismatch} uses an
independent wind-field realization with
$\eta_w=1.0$ and $\sigma_g=5$\,m\,s$^{-1}$, giving a disturbance
strength of $5.0$\,m\,s$^{-1}$, or approximately $5.6$ times the
nominal value. The combined condition applies four perturbations
simultaneously: mass $+20\%$, thrust efficiency $0.85$, a one-step
actuator delay, and this independent wind shift. It does not include
the inertia perturbation or the two-step delay.

\begin{figure}[tb]
	\centering
	\includegraphics[width=0.92\linewidth]{figures/fig06_robustness.pdf}
	\caption{Conflict-rate sensitivity to plant--model mismatch.}
	\label{fig:mismatch-robustness}
\end{figure}

The CBF-constrained configurations remain below
\textsc{Vanilla-MPC} in every tested regime. Taking \method{} as the
CBF reference, the conflict-rate gap to \textsc{Vanilla-MPC} ranges
from $10.5$ to $64.5$ percentage points. The gap narrows to
$10.5$ points under the reduced-mass condition, where the lighter
plant increases the available response margin, but widens under
several adverse perturbations. It reaches $57.5$ points for
mass $+20\%$, $52.0$ points under reduced thrust efficiency,
$59.5$ points under the two-step actuator delay, and $64.5$ points
under the combined four-factor shift.

Among the single-factor perturbations, the two-step actuator delay
produces the highest conflict rate for the certificate-free
configuration, $87.5\%$, and also raises the conflict rate of
\method{} from $11.0\%$ to $28.0\%$. This sensitivity is consistent
with the manoeuvring-authority diagnosis: delay reduces the rate at
which corrective commands can be realized, even when the nominal
acceleration authority available to the planner is unchanged.

A separate wind-scaling sweep increases the disturbance strength from
$0.9$ to $4.5$\,m\,s$^{-1}$ on the same encounter stream; conflict
rates change by at most $1.5$ percentage points and mean minimum
separation by at most $0.05$\,m in any configuration, so this sweep is
treated as a robustness check on the nominal wind setting.

The mismatch experiments therefore do not suggest that the nominal
result is an artefact of exact plant--planner agreement. The magnitude
of the CBF advantage varies with the perturbation and can narrow when
the realized plant becomes easier to manoeuvre, as in the reduced-mass
case, but the advantage remains positive across all tested conditions.
Together with the feasibility, oracle, and acceleration-authority
analyses, these results identify two distinct layers of the closed-loop
mechanism: the separation constraint strongly affects how the
available manoeuvring capability is used, while the manoeuvring
envelope itself limits what the closed loop can achieve once the
encounter becomes sufficiently demanding.

\subsection{Operational and prediction-space effects of task-aligned training}
\label{sec:results-signal}

The preceding analyses show that, under the deployment configuration,
closed-loop conflict outcomes are strongly shaped by the planning
regime and the available manoeuvring authority. This does not by
itself determine whether planner-informed predictor training provides
additional value once the predictor has been exposed to the same
conflict encounters. We therefore compare Stage~2 with
\textsc{Stage-1b}, the displacement-fine-tuning baseline that starts
from the same Stage~1 checkpoint and is matched in encounter exposure,
training order, optimizer setup, and update budget. The comparison
tests whether the task-aligned training objective provides an
operational increment beyond displacement-only adaptation to the same
encounter distribution.

\paragraph{Matched operational comparison.}

\begin{figure}[tb]
	\centering
	\includegraphics[width=0.9\linewidth]{figures/fig07_matched.pdf}
	\caption{Matched operational differences between Stage~2 and Stage-1b.}
	\label{fig:matched-forest}
\end{figure}

Under the nominal deployment condition, \textsc{Stage-1b} achieves a
substantially lower open-loop ADE than Stage~2
($1.84$ versus $4.32$\,m), but the two configurations produce nearly
identical closed-loop outcomes. Their conflict rates are $11.5\%$ and
$11.0\%$, respectively, with two encounters conflicting only under
\textsc{Stage-1b} and one only under Stage~2
(exact McNemar $p=1.00$). The paired difference in mean minimum
separation, defined as Stage~2 minus \textsc{Stage-1b}, is
$-0.04$\,m with a $95\%$ confidence interval of
$[-0.35,\,0.27]$\,m. The corresponding control-effort difference is
$-0.55$ with a $95\%$ interval of $[-1.68,\,0.58]$.

Maximum lateral offset is included as a supplementary
trajectory-deviation diagnostic rather than as the deviation term
used in Stage~2 training. Its paired difference is $+0.52$\,m
($95\%$ CI $[-1.11,\,2.14]$\,m), again with no detected difference
between the two configurations. Figure~\ref{fig:matched-forest}
summarizes these paired effects. For the conflict-rate comparisons in
panel~(a), the counts printed at the right give the discordant
episodes as
(\textsc{Stage-1b}-only conflict, Stage~2-only conflict).

The nominal result is not reversed under the selected plant--model
mismatch conditions. Figure~\ref{fig:matched-forest}(a) additionally
compares the same pair under mass $+20\%$, thrust efficiency $0.85$,
a two-step actuator delay, and the combined four-factor shift defined
in the preceding robustness analysis. The corresponding Stage-1b
conflict rates are $17.0\%$, $16.0\%$, $30.0\%$, and $31.0\%$,
compared with $17.5\%$, $16.5\%$, $28.0\%$, and $30.0\%$ for
Stage~2. All paired risk-difference intervals include zero. Thus, the
absence of a detectable operational increment in the nominal
comparison is not reversed by these selected model-mismatch
conditions.

We next examine sensitivity to Stage~2 training. The primary Stage~2
checkpoint is a separate independent training run, distinct from the
three additional robustness checkpoints. The three additional
checkpoints, trained independently under the same training-time
planner, yield deployment conflict rates of $12.0\%$, $11.0\%$, and
$11.5\%$, compared with $11.0\%$ for the primary checkpoint. Their
corresponding ADEs are $7.93$, $5.77$, and $7.16$\,m, compared with
$4.32$\,m for the primary checkpoint. The additional runs have a mean
deployment conflict rate of $11.5\pm0.5\%$ and mean ADE of
$6.96\pm1.09$\,m. Each run is compared with the fixed Stage-1b
baseline on the paired encounter outcomes, and none shows a detected
conflict-rate difference (exact McNemar $p=1.00$ in all three
comparisons).

The primary checkpoint therefore lies below the replicate range in
open-loop ADE while remaining within the replicate range in
deployment conflict rate, a descriptive observation consistent with
the open-loop/closed-loop decoupling found in the main comparison.

A further Stage~2 checkpoint trained under the deployment planner has
an ADE of $9.05$\,m and produces a deployment conflict rate of
$13.5\%$. Across all five Stage~2 checkpoints, deployment conflict
rates occupy a narrow $11.0$--$13.5\%$ range despite substantial
variation in training run and open-loop error. The three seed
comparisons share the same Stage-1b outcome vector and are therefore
used as training-stochasticity robustness checks rather than as
independent tests of seed-to-seed differences. The same checkpoints
are cross-evaluated under the training-time planner later in the
analysis.



Taken together, the matched comparison and the robustness checks show
no detected operational increment from the task-aligned objective
under the deployment planning configuration. This does not imply that
the training procedures produce equivalent predictors: their
open-loop errors differ, and the directional analysis below identifies
a structured prediction-space difference. The cross-configuration
analysis later tests whether these predictor-side differences retain
the same operational expression when the downstream planner changes.

\paragraph{Directional prediction-space signature.}

The task-aligned loss can influence the predictor through the planning
layer only where the planned actions are sensitive to predicted
neighbour positions. Because per-step active-set logs are not
available, we use a closest-point-of-approach-centred window,
$|k-k_{\mathrm{CPA}}|\leq3$, as an operational proxy for the region
in which separation constraints are most likely to be consequential.
Prediction errors in this window are decomposed into a signed
line-of-sight component $e_{\parallel}$ and an orthogonal magnitude
$|e_{\perp}|$. Negative $e_{\parallel}$ means that the neighbouring
aircraft is predicted closer to the ego aircraft than it is in the
realized trajectory. Because multiple prediction steps from the same
encounter are correlated, inferential comparisons are performed on
episode-level aggregates.

\begin{table}[tb]
	\centering
	\small
	\caption{Directional decomposition of prediction error in the CPA-centred window.}
	\label{tab:errdir}
	\begin{tabular}{lccc}
		\toprule
		Predictor
		& $e_{\parallel}$ (m, mean $\pm$ SE)
		& frac.\ toward ego
		& $|e_{\perp}|$ (m) \\
		\midrule
		Stage 1 (context)
		& $-0.70 \pm 0.49$
		& $53.8\%$
		& 8.38 \\
		Stage-1b
		& $-0.18 \pm 0.24$
		& $45.4\%$
		& 0.68 \\
		Stage 2 (\method)
		& $-0.34 \pm 0.24$
		& $54.6\%$
		& 2.91 \\
		\midrule
		Stage 2 $-$ Stage-1b
		& \multicolumn{3}{l}{
		$-0.164$\,m;
		$95\%$ CI $[-0.283,-0.045]$;
		paired $t$ $p=0.007$;
		Wilcoxon $p=0.001$} \\
		\bottomrule
	\end{tabular}
\end{table}

\begin{figure}[tb]
	\centering
	\includegraphics[width=0.95\linewidth]{figures/fig09_fingerprint.pdf}
	\caption{Directional prediction-error differences between Stage~2 and Stage-1b.}
	\label{fig:fingerprint}
\end{figure}

The configuration-specific means in Table~\ref{tab:errdir} should be
read descriptively. None of the three individual means provides the
appropriate test of the training objective, because the encounter
geometry is shared across configurations. The relevant quantity is
the paired Stage~2--Stage-1b difference. Stage~2 shifts the
episode-mean line-of-sight error by $-0.164$\,m relative to
displacement-only fine-tuning
($95\%$ paired-$t$ CI $[-0.283,-0.045]$\,m, $p=0.007$).
A Wilcoxon signed-rank test gives the same qualitative result
($p=0.001$). The bootstrap interval displayed in
Figure~\ref{fig:fingerprint}, $[-0.280,-0.046]$\,m, closely agrees
with the paired-$t$ interval.

The direction of this shift is toward the ego aircraft along the
line of sight and is therefore locally conservative with respect to
the separation geometry. The magnitude is small in operational scale:
$0.164$\,m is approximately $0.5\%$ of the $30$\,m separation
standard. At the encounter level, $120$ of the $200$ encounters have
a more negative episode-mean $e_{\parallel}$ under Stage~2 than under
\textsc{Stage-1b}. The orthogonal error magnitude is descriptively
larger for Stage~2 ($2.91$ versus $0.68$\,m), but no paired inference
is attached to this quantity. The evidence therefore supports a small
directional change in prediction space rather than a general
improvement in prediction accuracy.

Encounter~93 is an extreme observation under all three learned
predictors and strongly affects their configuration-specific means.
It does not drive the matched result. Excluding this encounter changes
the Stage~2--Stage-1b paired difference only from $-0.164$ to
$-0.171$\,m, with the corresponding $95\%$ interval remaining below
zero ($[-0.290,-0.052]$\,m; $p=0.005$). The paired directional effect
is therefore substantially less sensitive to this encounter than the
individual predictor means.

Stage~1 is retained in Table~\ref{tab:errdir} only as contextual
information. It differs from Stage~2 in both encounter exposure and
prediction-error scale and therefore does not provide the matched
comparison required to attribute the directional difference to the
fine-tuning objective.

\paragraph{Accuracy allocation around closest approach.}

A directional shift does not establish the stronger hypothesis that
task-aligned training reallocates predictive accuracy toward the
region in which separation decisions are made. To examine that
hypothesis, we compare prediction error inside the CPA-centred window
with error outside it. The CPA window contains all samples satisfying
$|k-k_{\mathrm{CPA}}|\leq3$; the outside-window aggregate contains
all available samples with $|k-k_{\mathrm{CPA}}|>3$.

\begin{table}[tb]
	\centering
	\small
	\caption{Prediction error inside and outside the CPA-centred window.}
	\label{tab:rq5-profile}
	\begin{tabular}{lccc}
		\toprule
		Predictor
		& CPA window (m)
		& Outside window (m)
		& window/outside \\
		\midrule
		Stage 1
		& 13.53 & 23.14 & 0.585 \\
		Stage-1b
		& 1.04 & 2.08 & 0.499 \\
		Stage 2 (\method)
		& 4.07 & 4.40 & 0.925 \\
		\bottomrule
	\end{tabular}
\end{table}

\begin{figure}[tb]
	\centering
	\begin{subfigure}{\linewidth}
		\includegraphics[width=\linewidth]
		{figures/fig08a_profile_stage1.pdf}
		\caption{Stage~1 prediction-error profile.}
		\label{fig:profile-a}
	\end{subfigure}\\[4pt]
	\begin{subfigure}{\linewidth}
		\includegraphics[width=\linewidth]
		{figures/fig08b_profile_matched.pdf}
		\caption{Stage-1b and Stage~2 prediction-error profiles.}
		\label{fig:profile-b}
	\end{subfigure}
	\caption{Prediction-error profiles around closest approach.}
	\label{fig:error-profile}
\end{figure}

Figure~\ref{fig:error-profile} shows the local error profiles through
seven steps from closest approach. The displayed range is intentionally
limited to $|k-k_{\mathrm{CPA}}|\leq7$ to resolve the behaviour around
the CPA-centred window; the outside-window values in
Table~\ref{tab:rq5-profile} are computed from all available prediction
samples satisfying $|k-k_{\mathrm{CPA}}|>3$, including samples beyond
the range displayed in the figure. The figure and table therefore
serve complementary purposes: the former shows the local profile,
whereas the latter reports the full outside-window aggregate.

The hypothesized concentration of prediction accuracy toward the
CPA-centred region is not supported under either an absolute or a
relative reading. In absolute terms, Stage~2 has a larger mean error
inside the CPA window than the displacement-fine-tuned Stage-1b
baseline ($4.07$ versus $1.04$\,m). Relative to the error outside the
window, the Stage~2 ratio is also larger
($0.925$ versus $0.499$). Thus, task-aligned training does not reduce
prediction error preferentially in the CPA-centred proxy region
relative to matched displacement-only fine-tuning.

The prediction-space analyses therefore refine the operational null
rather than overturn it. Under the deployment CBF configuration, the
task-aligned objective produces a small, statistically detectable
line-of-sight shift in the CPA-centred window, but it neither improves
absolute prediction accuracy in that region nor concentrates accuracy
there relative to the outside-window error. These changes have no
detectable counterpart in conflict rate, minimum separation, control
effort, or the supplementary lateral-offset diagnostic under the
matched deployment comparison. The conclusion is specific to this
constraint-governed regime; the following analyses examine the
operating conditions under which the task-aligned objective regains
closed-loop leverage.


\subsection{Boundary conditions of the closed-loop findings}
\label{sec:results-boundary}

The preceding analyses were conducted at a specific operating point:
the deployment planning configuration and a short tactical response
window. The conclusions drawn at that point need not hold when either
the time available for conflict resolution or the downstream planning
configuration changes. We therefore examine two boundary conditions.
The first varies the available detection lead time to determine whether
the residual conflicts attributed to limited manoeuvring authority
persist when the aircraft is given more time to respond. The second
cross-evaluates Stage~2 checkpoints under the training-time and
deployment planning configurations to determine whether predictor
differences retain the same closed-loop expression when the downstream
planner changes.

\paragraph{Boundary in available detection lead time.}

The residual-conflict analysis showed that replacing the learned
forecast with the true future trajectory did not resolve the conflicts
remaining under the deployment configuration. One possible
explanation is that the conflict is identified too late for either a
learned or an oracle predictor to make a material difference. To test
this possibility, we vary the available time to closest approach while
holding the encounter construction and deployment planner otherwise
fixed. The lead-time sweep uses separately sampled encounters generated from
the same held-out trajectory pool and the same encounter-construction
procedure as the primary evaluation. At the $2$\,s condition, the
nominal closest-approach timing matches that of the primary evaluation,
but the sampled episodes are not identical. Results are therefore
compared at the aggregate level, and no episode-wise paired inference
is made between the primary and lead-time evaluation sets. 

\begin{table}[tb]
	\centering
	\small
	\caption{Conflict rate across available detection lead times.}
	\label{tab:leadtime}
	\begin{tabular}{lcccc}
		\toprule
		Lead time
		& Stage 1
		& Stage-1b
		& \method
		& Oracle \\
		\midrule
		$1$\,s  & 86.0 & 83.0 & 82.5 & 82.0 \\
		$2$\,s  & 13.0 & 11.5 & 11.0 & 11.0 \\
		$3$\,s  &  1.5 &  0.0 &  0.0 &  0.0 \\
		$4$\,s  &  7.0 &  0.0 &  0.0 &  0.0 \\
		$5$\,s  & 10.0 &  0.5 &  0.0 &  0.0 \\
		$6$\,s  & 12.0 &  0.5 &  0.0 &  0.0 \\
		$7$\,s  &  8.5 &  0.5 &  0.0 &  0.0 \\
		$10$\,s &  0.5 &  0.0 &  0.0 &  0.0 \\
		$20$\,s &  0.0 &  0.0 &  0.0 &  0.0 \\
		\bottomrule
	\end{tabular}
\end{table}

\begin{figure}[tb]
	\centering
	\includegraphics[width=0.98\linewidth]{figures/fig10_leadtime.pdf}
	\caption{Conflict rate across available detection lead times.}
	\label{fig:leadtime}
\end{figure}

The effect of additional response time is sharp at the tactical
endpoint. At a lead time of $1$\,s, conflict rates remain above
$82\%$ for Stage-1b, Stage~2, and the oracle reference. At $2$\,s,
the corresponding rates fall to $11.5\%$, $11.0\%$, and $11.0\%$.
No Stage-1b, Stage~2, or oracle conflicts are observed at the
$3$\,s point, and Stage~2 and the oracle remain conflict-free at all
longer lead times on the tested grid.

The episode-level attribution is consistent with the mechanism
identified above. Of the $165$ Stage~2 conflicts at $1$\,s,
approximately $98\%$ remain unresolved under the corresponding oracle
reference and are therefore classified as
manoeuvring-authority-limited under the fixed planning and control
envelope. At $2$\,s, the same classification applies to approximately
$95\%$ of the $22$ Stage~2 conflicts arising in that lead-time
condition. From $3$\,s onward Stage~2
has no residual conflicts to attribute. The comparison therefore
indicates that additional response time can remove conflicts for which
improved prediction alone provides little leverage at shorter lead
times.

The transition between the tested $2$- and $3$-s points initially
raises a potential confounding issue because the deployment planning
horizon is also $3.0$\,s
($H_p=15$ at $\Delta t=0.2$\,s). We therefore repeat the Stage~2
lead-time sweep with planning horizons of $H_p=8$, $15$, and $25$,
corresponding to $1.6$, $3.0$, and $5.0$\,s.

\begin{table}[tb]
	\centering
	\small
	\caption{Stage~2 lead-time transition across planning horizons.}
	\label{tab:leadtime-horizon}
	\begin{tabular}{lccc}
		\toprule
		Lead time
		& $H_p=8$ (1.6\,s)
		& $H_p=15$ (3.0\,s)
		& $H_p=25$ (5.0\,s) \\
		\midrule
		$1$\,s & 82.0 & 82.5 & 82.5 \\
		$2$\,s & 11.0 & 11.0 & 11.0 \\
		$3$\,s &  0.0 &  0.0 &  0.0 \\
		\bottomrule
	\end{tabular}
\end{table}

The transition is unchanged across the three planning horizons. All
three produce an $11.0\%$ conflict rate at $2$\,s and no observed
conflicts at $3$\,s; the same zero-conflict outcome continues at
$4$, $5$, and $7$\,s. In particular, the transition remains at the
same lead time when the planning horizon is substantially shorter
($1.6$\,s) or longer ($5.0$\,s) than $3$\,s. The coincidence
between the deployment planning horizon and the observed $3$\,s
transition is therefore not sufficient to explain the result within
the tested horizon range.

The numerical value of this transition should nevertheless not be
interpreted as a universal response-time requirement. It is specific
to the encounter geometries, closing-speed distribution, $30$\,m
separation standard, and manoeuvring envelope used here. The result
instead establishes a more limited operational point: for the tested
encounters, increasing the available response window from $2$ to
$3$\,s eliminates conflicts that persist even when the planner is
given the true future trajectory, and this result is robust to planning
horizons from $1.6$ to $5.0$\,s.

Stage~1 follows a different pattern. It has the highest conflict rate
of the four predictors at both tactical endpoints and develops a
non-monotone response at intermediate lead times, rising from
$1.5\%$ at $3$\,s to $12.0\%$ at $6$\,s before falling again.
This intermediate increase is absent from Stage-1b, Stage~2, and the
oracle reference. Repeating the experiment across the three planning
horizons preserves the same qualitative Stage~1 pattern, while the
oracle remains conflict-free over the corresponding intermediate
range. The non-monotonicity is therefore not explained by an
interaction between lead time and planning-horizon length. It is
consistent with a predictor-specific effect associated with the
Stage~1 forecast under the conflict-encounter distribution, although
the present analysis does not identify its trajectory-level cause.

\paragraph{Planner-dependent expression of predictor differences.}

The matched comparison in the preceding analysis found no detectable
operational increment from the task-aligned objective under the
deployment planning configuration. This deployment result, however,
does not establish that predictor-side differences would remain
similarly muted under another downstream planner. We therefore
cross-evaluate the Stage~2 checkpoints under two planning
configurations: the deployment configuration
\[
(\gamma,H_p,a_{\max})=(0.1,15,20\,\mathrm{m\,s^{-2}})
\]
and the training-time configuration
\[
(\gamma,H_p,a_{\max})=(0.4,8,10\,\mathrm{m\,s^{-2}}).
\]
The latter is the planner used to train the primary Stage~2 checkpoint
and its three independently trained robustness replicates. Because the
two planning configurations differ jointly in barrier coefficient,
planning horizon, and acceleration authority, the comparison is not
interpreted as a one-factor test of constraint strength.

\begin{figure}[tb]
	\centering
	\includegraphics[width=0.96\linewidth]
	{figures/fig12_planner_expression.pdf}
	\caption{Planner-dependent expression of Stage~2 predictor differences.}
	\label{fig:planner-expression}
\end{figure}

Figure~\ref{fig:planner-expression} shows a pronounced interaction
between predictor checkpoint and evaluation planner. Under the
deployment planner, the five Stage~2 checkpoints occupy a narrow
conflict-rate range of $11.0$--$13.5\%$, a span of $2.5$ percentage
points. The four checkpoints trained under the training-time planner
produce rates of $11.0\%$, $12.0\%$, $11.0\%$, and $11.5\%$ in this
deployment regime, while the checkpoint trained under the deployment
planner yields $13.5\%$.

The same checkpoints separate sharply under the training-time planner.
The primary Stage~2 checkpoint and its three independently trained
replicates produce conflict rates of $28.5\%$, $27.5\%$, $28.5\%$,
and $27.0\%$, respectively. Their mean is $27.9\pm0.8\%$, and the
episode-level improvement relative to Stage-1b is nested in all four
runs: each resolves between $55$ and $58$ Stage-1b conflict episodes
while introducing none. For the primary checkpoint, Stage-1b conflicts
in $56.0\%$ of the encounters and Stage~2 in $28.5\%$, with
$55$ against $0$ discordant episodes
(exact McNemar $p=5.6\times10^{-17}$).

The Stage~2 checkpoint trained under the deployment planner behaves
very differently under the same training-time evaluation planner. Its
conflict rate rises to $69.0\%$, $13.0$ percentage points above the
Stage-1b reference. The paired outcome pattern is no longer nested:
$9$ encounters conflict only under Stage-1b and $35$ only under this
Stage~2 checkpoint (exact McNemar
$p=1.06\times10^{-4}$). Relative to the primary Stage~2 checkpoint
under the same planner, the difference is larger still: all conflicts
of the primary checkpoint are contained in the deployment-trained
checkpoint's conflict set, which adds $81$ further conflict episodes
($0$ against $81$ discordant episodes; exact McNemar
$p=8.27\times10^{-25}$).

Average displacement error does not provide a monotone explanation for
this divergence. The Stage-1b, primary Stage~2, and deployment-trained
Stage~2 predictors have ADEs of $1.84$, $4.32$, and $9.05$\,m,
respectively, while Stage~1 has an ADE of $20.90$\,m. Under the
training-time planner, their conflict rates are $56.0\%$, $28.5\%$,
$69.0\%$, and $56.5\%$, respectively. The large variation in
closed-loop outcome therefore cannot be read from ADE ordering alone.

The result is not evidence for a general matching rule in which a
predictor performs best whenever its training and evaluation planners
coincide. The primary Stage~2 checkpoint, for example, has a lower
absolute conflict rate under the non-matching deployment planner than
under its own training-time planner. Nor does the present experiment
identify which change in $\gamma$, $H_p$, or $a_{\max}$ produces the
modulation. What the cross-evaluation establishes is narrower: the
downstream planning configuration strongly affects how differences
among predictor checkpoints are expressed in closed-loop conflict
outcomes. Those differences are compressed under the deployment
planner but become much larger under the training-time planner.

The four independently trained checkpoints on the training-time side
show that the favourable $27$--$28.5\%$ result is reproducible across
training stochasticity, including the nested episode pattern. Only one
Stage~2 checkpoint is available from training under the deployment
planner, however, so training-run stochasticity cannot be fully
separated from training-configuration effects on that side.

Together, the two boundary analyses qualify the earlier findings in
different ways. The residual conflicts observed at short lead times
can be removed by increasing the available response window, even
without improving prediction accuracy. Separately, predictor
differences that are weakly expressed under the deployment planner can
become large under another planning configuration. The numerical
transition times and cross-configuration effect sizes are specific to
the tested encounter and aircraft settings. More broadly, prediction
quality and predictor training cannot be assigned a fixed operational
value independently of the downstream planning and manoeuvring
environment in which they are evaluated.



\subsection{System-level effects of mixed avoidance logics}
\label{sec:results-corridor}

The preceding experiments use controlled conflict encounters to
isolate the effects of prediction, planning, training, manoeuvring
authority, and response time. Such encounter-level comparisons do not
capture a different operational question that arises under partial
deployment: how aircraft governed by different avoidance logics
interact when they share the same traffic stream. We therefore conduct
an exploratory corridor-scale experiment in which
\method-equipped aircraft operate together with aircraft controlled by
optimal reciprocal collision avoidance (ORCA)
\citep{vandenberg2011reciprocal}. The purpose of this experiment is
not to rank the two avoidance methods in isolation, but to examine how
system-level outcomes and group-specific performance change as the
fleet composition varies.

\paragraph{Corridor setting and mixed-equipage design.}

The experiment uses a bidirectional corridor with aircraft entering
from opposite ends under discrete-time Bernoulli arrivals. The
corridor is $600$\,m long with a half-width of $120$\,m at an
operating altitude of $100$\,m. Aircraft cruise at $25$\,m\,s$^{-1}$
and interact with up to their three nearest neighbours within a
$250$\,m sensing radius. \method-equipped aircraft use the deployment
prediction and CBF-constrained planning stack, whereas unequipped
aircraft use a fixed ORCA reciprocal-avoidance controller. The
equipped fraction is varied over
$p\in\{0,25,50,75,100\}\%$.

Two demand levels are considered, with per-step arrival probabilities
of $0.16$ and $0.06$ at each end of the corridor. Each
demand--penetration condition is evaluated over six independent
replications, with a $100$-step warm-up followed by a $400$-step
evaluation horizon. Throughput is the number of completed corridor
transits per minute. Delay is measured among completed transits as the
excess travel time relative to free flow. An aircraft is counted as a
conflict agent if its realized separation falls below the operational
standard $d_{\mathrm{sep}}=30$\,m at least once during its transit.
Aircraft leaving the prescribed lateral operating region are recorded
separately as lateral exits. No timeout removals occur in these
experiments, so non-completion is attributable to lateral exit.

Both avoidance logics are evaluated against the same $30$\,m
operational separation standard, but their internal constraint
mechanisms are not dynamically equivalent. In particular, the
CBF-constrained planner regulates the evolution of separation in
addition to the geometric threshold, whereas ORCA uses reciprocal
velocity constraints. The corridor experiment should therefore be
read as a mixed-logic interaction study under the specified
parameterizations rather than as a matched comparison of two planners.

\begin{figure}[tb]
	\centering
	\includegraphics[width=\linewidth]{figures/fig13_corridor.pdf}
	\caption{Corridor-scale outcomes across mixed-equipage levels.}
	\label{fig:corridor-penetration}
\end{figure}

\paragraph{Throughput and delay.}

Figure~\ref{fig:corridor-penetration} summarizes the corridor
outcomes at both demand levels. At high demand, mean system throughput
does not decrease as the equipped fraction grows. It changes from
$44.8\pm2.1$ passes/min at $p=0\%$ to
$47.7\pm5.2$ passes/min at $p=100\%$, with intermediate values of
$44.8$, $45.3$, and $46.8$ passes/min. The magnitude of this change
is comparable to the between-replication variability, and we
therefore interpret throughput as broadly maintained rather than as
evidence of a systematic capacity increase.

The clearer system-level pattern appears in completed-trip delay.
Under high demand, mean system delay decreases from
$6.5\pm1.4$\,s at $p=0\%$ to $2.6\pm0.6$\,s at
$p=100\%$, with intermediate means of $4.3$, $3.9$, and
$3.0$\,s. The ORCA-controlled subgroup shows the same direction over
the mixed-equipage range: although its controller is unchanged, its
mean delay decreases from $6.5$\,s at $p=0\%$ to
$5.6$, $5.5$, and $4.4$\,s at $p=25\%$, $50\%$, and
$75\%$, respectively. No ORCA subgroup exists at full equipage.

Because the ORCA control law itself is held fixed across penetration
levels, this decline provides an indicator of a positive cross-group
efficiency effect: unequipped aircraft experience lower delay as a
larger share of the surrounding traffic uses the equipped planning
stack. This comparison is not paired across penetration levels,
however, because the traffic realizations and group composition also
change. It is therefore interpreted as a system-level externality
pattern rather than as an isolated causal effect of equipage on an
otherwise identical ORCA trajectory.

The lower-demand experiment provides a second loading condition for
the same interaction. Under low demand, system throughput is
non-monotone, reaching its lowest value near $75\%$ equipage before
partially recovering at full equipage. Meanwhile, the endpoint system
delay decreases from $2.9\pm0.9$ to $1.4\pm0.5$\,s between zero
and full equipage. The intermediate delay values are less monotone
than in the high-demand case, but the endpoint comparison is
consistent with the same qualitative reduction in completed-trip
delay.

\paragraph{Conflict outcomes and cross-group interaction.}

The conflict-rate pattern is less monotone than the delay response.
At high demand, the system-level conflict rate changes from
$50.3\pm5.6\%$ at zero equipage to
$57.2\pm3.5\%$, $53.3\pm6.7\%$, and
$52.1\pm5.4\%$ at intermediate penetration levels, before reaching
$47.7\pm12.2\%$ at full equipage. The increase at intermediate
penetration is modest relative to the observed between-replication
variability and is therefore treated descriptively rather than as an
established deterioration in safety.

The hollow markers in the conflict-rate panels of
Figure~\ref{fig:corridor-penetration} show the six
replication-level outcomes for the ORCA-controlled subgroup at each
penetration level for which that subgroup is present. Their dispersion
illustrates the variability underlying the aggregate means. A
qualitatively non-monotone pattern also appears under the lower-demand
condition. Taken together, these results do not establish a monotonic
safety benefit from increasing equipage. They are instead consistent
with the possibility that mixed avoidance logics introduce
coordination effects at intermediate penetration, a mechanism that
would require targeted pair-type or trajectory-level analysis to
identify.

The corridor conflict rates are intentionally interpreted as
stress-test descriptors rather than estimates of operational
population risk. Their relatively high absolute levels also make the
experiment more informative about interaction patterns across control
logics than about the absolute safety performance of either avoidance
method.

\paragraph{Corridor retention and interpretation.}

The reduction in completed-trip delay is accompanied by a second,
less favourable trend. Under high demand, the lateral-exit share
increases from $0\%$ at zero equipage to approximately
$1.5\%$, $4.2\%$, $4.4\%$, and $8.0\%$ as the equipped
fraction rises. Under low demand, the lateral-exit share rises from approximately
$0.9\%$ at zero equipage to a peak of $10.8\%$ at $75\%$ equipage,
before decreasing slightly to approximately $9.9\%$ at full equipage. Because no
aircraft time out, the completion share is the complement of the
lateral-exit share and is not plotted separately.

This trend is important when interpreting the delay results. Delay is
computed only for aircraft that complete the corridor, whereas
lateral exits remain part of the broader system outcome. The decline
in completed-trip delay therefore coexists with a growing
corridor-retention cost as equipage increases. The present experiment
does not isolate the trajectory-level cause of these exits, although
the pattern is consistent with stronger lateral avoidance becoming
more consequential within a bounded corridor.

Overall, the corridor experiment reveals a trade-off rather than a
uniform benefit from increasing equipage. System throughput is broadly
maintained, and completed-trip delay decreases, including for the
unequipped ORCA-controlled subgroup under high demand. These efficiency
patterns coexist with an increasing lateral-exit share and with
conflict rates that vary non-monotonically across mixed-equipage
levels. The results therefore suggest that partial deployment should
be evaluated at the traffic-system level, where interactions between
different avoidance conventions can generate externalities that are
not visible in isolated encounter comparisons.





\FloatBarrier
\section{Conclusions}
\label{sec:conclusions}

This study set out to test two premises that govern how prediction is
built and judged for tactical separation in low-altitude corridors:
that a more accurate forecast buys a safer operation, and that letting
the planner supervise the predictor buys more still. In a closed loop
where separation is enforced as an explicit constraint, neither premise
held in the form in which it is normally assumed. The reason is
structural rather than incidental. A constraint layer that is doing its
job absorbs forecast differences before they reach the trajectory, so
the very mechanism that makes the operation safe also makes it
insensitive to the quantity by which predictors are ranked. Prediction
error and operational risk are not two measurements of one underlying
quality; they are separated by a layer whose transfer characteristics
are themselves a design variable.

That observation reorganizes where safety argumentation should be
directed. In the configuration studied here, the binding quantities
were the manoeuvring authority available to execute an avoidance and
the time available to begin it, not the fidelity of the forecast that
triggered it. Both are properties of the vehicle and of the airspace
rather than of the learned model, and both are already instruments of
transportation policy. This suggests that corridor admission criteria
should carry an explicit manoeuvre-performance floor, and that
surveillance and corridor geometry should be evaluated by the response
window they deliver to the aircraft that must use it. It also suggests
a specific caution for the certification of learned components: a
predictor validated only by displacement error, or validated in the
loop of one planner configuration, has not been shown to be safe in
the loop of another. Our cross-evaluation produced exactly that
failure, with predictor differences compressed under one planner and
reversed in sign under a second.

For research practice, the implication is that a prediction model
intended for separation should be reported together with the decision
layer through which it will act, and with the envelope in which that
layer must operate. Displacement error remains a legitimate measure of
forecast quality; what does not follow is that improving it improves
operations. Where a decision-focused objective is used, our results
argue for reporting its effect under more than one downstream
configuration, since the operational value we measured was not a stable
attribute of the trained predictor. At the system level, the corridor
experiment indicated that partial deployment of a new avoidance logic
can help aircraft that have not adopted it, while shifting cost into a
different currency, here corridor retention rather than delay. Such
transfers between outcome dimensions are invisible in encounter-level
comparison and argue for evaluating deployment programmes at the
traffic-system level.

Several limitations bound these conclusions and indicate what would
sharpen them. The controlled experiments resolve one neighbouring
aircraft, so simultaneous binding constraints are represented only in
the exploratory corridor extension and never with feasibility
attribution. The evaluation encounters are stress-test constructions,
which makes the absolute rates unsuitable as estimates of operational
risk even though the paired contrasts remain interpretable. The
cross-planner comparison varies barrier coefficient, horizon, and
acceleration limit jointly, so it establishes that the downstream
configuration governs the expression of predictor differences without
identifying which factor is responsible; a single-factor sweep would
resolve this, and would also determine whether the effect is graded or
threshold-like. Only one predictor was trained under the deployment
configuration, leaving training stochasticity and configuration effects
entangled on that side of the comparison. The corridor results rest on
six replications of one geometry and are reported without significance
testing. Finally, the wind field, vehicle model, and separation
standard are fixed throughout, and the observed lead-time threshold is
specific to the closing-speed distribution of the encounters used here;
whether the same threshold governs other traffic mixes is an open
question, though its invariance across planning horizons suggests it is
set by encounter geometry rather than by planner design.

The broader conclusion we draw is a methodological one. Modular
autonomy is normally validated module by module, on the assumption that
component quality composes into system quality. The evidence assembled
here indicates that, for separation-critical operations, it does not
compose in any fixed way: the same predictor was operationally
indistinguishable from its baseline under one planner and substantially
worse under another, without its forecasts changing at all. Assigning
operational value to a prediction model therefore requires naming the
loop in which it will act. Building instruments that can replay an
encounter while exchanging one layer at a time, and reporting what such
instruments reveal even when the answer is that a favoured mechanism
did not help, seems to us a precondition for admitting learned
components into shared airspace.

%\bibliographystyle{elsarticle-num}
\bibliographystyle{plainnat}
\bibliography{refs}

\FloatBarrier
\appendix
\section{Proofs}
\label{app:proofs}

This appendix supplies the arguments deferred from
Section~\ref{sec:method}. Appendix~\ref{app:invariance} proves
Proposition~\ref{prop:forward-invariance}, on the clearance guaranteed
by a feasible barrier recursion. Appendix~\ref{app:gradient} proves
Proposition~\ref{prop:gradient}, on differentiability of the planning
solution map, after a preliminary lemma establishing that the map from
predicted positions to program data is itself smooth.
Appendix~\ref{app:support} proves
Proposition~\ref{prop:gradient-support}, which localizes the
operational gradient to active barrier inequalities. Throughout,
$t$ is a fixed planning instant, and we abbreviate the planned
quantities as in~\eqref{eq:planning-qp}, writing $\bm{p}_k$ for
$\bm{p}_{t+k\mid t}$ and $\hat{\bm{p}}^{\,i}_{k}$ for
$\hat{\bm{p}}^{\,i}_{t+k\mid t}$.

\subsection{Forward invariance of the planned clearance}
\label{app:invariance}

\begin{proof}[Proof of Proposition~\ref{prop:forward-invariance}]
	Fix a neighbour $i\in\mathcal{N}_t$ and suppose the
	program~\eqref{eq:planning-qp} admits a feasible point with
	$\bm{\epsilon}=\bm{0}$. Feasibility of that point requires, for
	every $k=1,\dots,H_p$,
	\begin{equation}
		\label{eq:app-recursion}
		h^{i}_{k}\;\geq\;(1-\gamma)\,h^{i}_{k-1},
		\qquad \gamma\in(0,1],
	\end{equation}
	where $h^i_k$ is the linearized clearance~\eqref{eq:cbf-barrier}
	evaluated on the planned trajectory and $h^i_0$ is evaluated at the
	current ego and neighbour positions under the convention
	$\hat{\bm{p}}^{\,i}_{t\mid t}:=\bm{p}^i_t$.

	\emph{Step 1: non-negativity along the horizon.} Since
	$\gamma\in(0,1]$ we have $1-\gamma\in[0,1)$, so the multiplier in
	\eqref{eq:app-recursion} is non-negative. Applying
	\eqref{eq:app-recursion} inductively from $k=1$ therefore preserves
	sign: if $h^i_{k-1}\geq0$ then
	$h^i_k\geq(1-\gamma)h^i_{k-1}\geq0$. With the hypothesis
	$h^i_0\geq0$ as the base case, induction on $k$ gives
	\begin{equation}
		\label{eq:app-geometric}
		h^{i}_{k}\;\geq\;(1-\gamma)^{k}\,h^{i}_{0}\;\geq\;0,
		\qquad k=0,1,\dots,H_p .
	\end{equation}
	The first inequality in~\eqref{eq:app-geometric} follows by chaining
	\eqref{eq:app-recursion} $k$ times, which is valid because each
	factor $(1-\gamma)$ is non-negative and the inequalities are
	therefore preserved under multiplication. Note that for $\gamma=1$
	the bound degenerates to $h^i_k\geq0$ for $k\geq1$, which is still
	the conclusion required; the geometric factor is informative only
	for $\gamma<1$.

	\emph{Step 2: from the linearized clearance to Euclidean
		separation.} Fix $k$ and let
	$\bm{n}^{i}_{k}$ be the unit linearization direction of
	Section~\ref{sec:method-planner}, so
	$\|\bm{n}^{i}_{k}\|_2=1$. By~\eqref{eq:cbf-barrier},
	\begin{equation}
		\label{eq:app-halfspace}
		h^{i}_{k}
		=\bm{n}^{i\,\top}_{k}
		\big(\bm{p}_{k}-\hat{\bm{p}}^{\,i}_{k}\big)
		-d_{\mathrm{sep}} .
	\end{equation}
	The Cauchy--Schwarz inequality gives
	$\bm{n}^{i\,\top}_{k}\bm{z}\leq\|\bm{n}^{i}_{k}\|_2\|\bm{z}\|_2
		=\|\bm{z}\|_2$ for any $\bm{z}\in\R^3$. Taking
	$\bm{z}=\bm{p}_{k}-\hat{\bm{p}}^{\,i}_{k}$ and
	using~\eqref{eq:app-halfspace},
	\begin{equation}
		\big\|\bm{p}_{k}-\hat{\bm{p}}^{\,i}_{k}\big\|_2
		\;\geq\;
		\bm{n}^{i\,\top}_{k}
		\big(\bm{p}_{k}-\hat{\bm{p}}^{\,i}_{k}\big)
		\;=\;
		h^{i}_{k}+d_{\mathrm{sep}}
		\;\geq\;
		d_{\mathrm{sep}},
	\end{equation}
	where the final inequality uses $h^i_k\geq0$
	from~\eqref{eq:app-geometric}. Since $k\leq H_p$ and
	$i\in\mathcal{N}_t$ were arbitrary, the planned trajectory keeps at
	least the required clearance from every predicted mean position at
	every planned step.
\end{proof}

\begin{remark}[What the argument does not deliver]
	\label{rem:invariance-scope}
	Three features of the proof delimit its reach, and each is measured
	rather than assumed in Section~\ref{sec:results-mechanism}.
	First, \eqref{eq:app-halfspace} compares the planned position with
	the \emph{predicted} neighbour position, so the conclusion transfers
	to the true traffic only insofar as the forecast is adequate at the
	planned passing geometry. Second, the argument is stated at one
	planning instant; because the planner is executed in receding
	horizon, the per-instant conclusions chain into a closed-loop
	property only if each replanned program again admits a feasible
	point with $\bm{\epsilon}=\bm{0}$ and inherits $h^i_0\geq0$, which
	is exactly what can fail when a refreshed prediction moves the
	keep-clear region faster than the acceleration envelope can answer.
	Third, the clearance lives on the double-integrator surrogate,
	whereas the aircraft flies~\eqref{eq:dynamics} through the saturated
	inner law, so surrogate--plant mismatch and the tilt and moment
	clamps bound how faithfully the planned motion is realized. The
	conservatism in Step~2 is one-sided and therefore harmless: the
	projection bound can only understate the true Euclidean separation.
\end{remark}

\subsection{Differentiability of the planning solution map}
\label{app:gradient}

Write the decision vector of~\eqref{eq:planning-qp} as
$\bm{z}=(\bm{a},\bm{p},\bm{v},\bm{\epsilon})\in\R^{n_z}$ and collect
the predicted mean positions available at instant $t$ into
$\bm{\mu}=\{\hat{\bm{p}}^{\,i}_{k}\}_{i\in\mathcal{N}_t,\,k\leq H_p}$.
The program is then
\begin{equation}
	\label{eq:app-qp}
	\min_{\bm{z}}\ \tfrac12\bm{z}^{\top}\!Q\,\bm{z}+\bm{c}^{\top}\bm{z}
	\quad\text{s.t.}\quad
	A\bm{z}=\bm{b},\qquad
	G(\bm{\mu})\,\bm{z}\;\leq\;\bm{g}(\bm{\mu}),
\end{equation}
where the equality block collects the initial conditions and the
double-integrator recursion, and the inequality block collects the
acceleration bounds, the slack sign conditions, and the barrier
recursion. Only the barrier rows depend on $\bm{\mu}$. The cost data
$(Q,\bm{c})$ are independent of $\bm{\mu}$, since the objective
of~\eqref{eq:planning-qp} involves only the reference route, the
accelerations, and the slacks.

A preliminary lemma is required, because $\bm{\mu}$ enters the barrier
rows twice: directly through the offset in~\eqref{eq:cbf-barrier}, and
through the linearization direction $\bm{n}^i_k$, which is itself a
normalization and hence non-smooth where its argument vanishes.

\begin{lemma}[Smoothness of the program data in the forecast]
	\label{lem:data-smooth}
	Let $\bm{\rho}^{i}_{k}
		:=\bm{r}_{t+k}-\hat{\bm{p}}^{\,i}_{k}$ denote the offset from the
	predicted position to its reference-route point, and suppose
	$\rho^{i}_{k}:=\|\bm{\rho}^{i}_{k}\|_2\geq\rho_{\min}>0$ for every
	$i$ and $k$. Then $\bm{n}^{i}_{k}
		=\bm{\rho}^{i}_{k}/\rho^{i}_{k}$ is $C^{\infty}$ in
	$\hat{\bm{p}}^{\,i}_{k}$ on that region, the maps
	$\bm{\mu}\mapsto G(\bm{\mu})$ and $\bm{\mu}\mapsto\bm{g}(\bm{\mu})$
	are $C^{\infty}$, and the derivative of the clearance
	\eqref{eq:cbf-barrier} with respect to the predicted position is
	\begin{equation}
		\label{eq:app-dh}
		\frac{\partial h^{i}_{k}}
		{\partial\hat{\bm{p}}^{\,i}_{k}}
		=-\Big(
		\bm{n}^{i}_{k}
		+\tfrac{1}{\rho^{i}_{k}}
		\big(I-\bm{n}^{i}_{k}\bm{n}^{i\,\top}_{k}\big)
		\big(\bm{p}_{k}-\hat{\bm{p}}^{\,i}_{k}\big)
		\Big)^{\!\top}.
	\end{equation}
\end{lemma}

\begin{proof}
	The map $\bm{y}\mapsto\bm{y}/\|\bm{y}\|_2$ is $C^{\infty}$ on
	$\R^3\setminus\{\bm{0}\}$, and $\hat{\bm{p}}^{\,i}_{k}\mapsto
		\bm{\rho}^{i}_{k}$ is affine, so the composition
	$\hat{\bm{p}}^{\,i}_{k}\mapsto\bm{n}^{i}_{k}$ is $C^{\infty}$
	wherever $\rho^{i}_{k}>0$, which the hypothesis
	$\rho^{i}_{k}\geq\rho_{\min}$ guarantees uniformly. Its Jacobian is
	the standard normalization derivative
	\begin{equation}
		\label{eq:app-dn}
		\frac{\partial\bm{n}^{i}_{k}}
		{\partial\hat{\bm{p}}^{\,i}_{k}}
		=-\frac{1}{\rho^{i}_{k}}
		\big(I-\bm{n}^{i}_{k}\bm{n}^{i\,\top}_{k}\big),
	\end{equation}
	the sign arising because $\bm{\rho}^{i}_{k}$ depends on
	$\hat{\bm{p}}^{\,i}_{k}$ with a minus sign. Differentiating
	\eqref{eq:cbf-barrier}, in which $\hat{\bm{p}}^{\,i}_{k}$ appears
	both in the subtracted position and inside $\bm{n}^{i}_{k}$, and
	applying the product rule,
	\begin{equation}
		\frac{\partial h^{i}_{k}}
		{\partial\hat{\bm{p}}^{\,i}_{k}}
		=-\bm{n}^{i\,\top}_{k}
		+\big(\bm{p}_{k}-\hat{\bm{p}}^{\,i}_{k}\big)^{\!\top}
		\frac{\partial\bm{n}^{i}_{k}}
		{\partial\hat{\bm{p}}^{\,i}_{k}} ,
	\end{equation}
	and substituting~\eqref{eq:app-dn} and using the symmetry of the
	projector $I-\bm{n}^{i}_{k}\bm{n}^{i\,\top}_{k}$
	gives~\eqref{eq:app-dh}. Each barrier row
	of~\eqref{eq:app-qp} is an affine function of the clearances
	$h^{i}_{k}$ with constant coefficients determined by $\gamma$, which
	is a compiled constant and not a program parameter, so smoothness of
	the clearances in $\bm{\mu}$ transfers to $G$ and $\bm{g}$.
\end{proof}

\begin{remark}[Role of $\rho_{\min}$]
	\label{rem:rhomin}
	The transverse term in~\eqref{eq:app-dh} carries the factor
	$1/\rho^{i}_{k}$ and therefore diverges as the predicted position
	approaches its reference-route point, where the linearization
	direction ceases to be defined. The hypothesis
	$\rho^{i}_{k}\geq\rho_{\min}$ in
	Proposition~\ref{prop:gradient} excludes precisely that degeneracy,
	which is why it appears alongside the constraint-qualification
	conditions rather than as a technical afterthought. In the
	encounters used in this study the predicted neighbour positions stay
	well away from the ego reference route, so the condition is not
	binding in practice; it is stated because the gradient expression is
	otherwise meaningless.
\end{remark}

\begin{proof}[Proof of Proposition~\ref{prop:gradient}]
	Let $\bar{\bm{\mu}}$ denote the current forecast and
	$(\bar{\bm{z}},\bar{\bm{\lambda}},\bar{\bm{\nu}})$ a primal--dual
	solution of~\eqref{eq:app-qp} at $\bar{\bm{\mu}}$, with
	$\bm{\lambda}$ the multipliers of the inequality block and
	$\bm{\nu}$ those of the equality block. Let
	$\mathcal{A}=\{j:\,\bm{G}_j(\bar{\bm{\mu}})\bar{\bm{z}}
		=g_j(\bar{\bm{\mu}})\}$ be the active set at
	$\bar{\bm{z}}$, and write $G_{\mathcal{A}},\bm{g}_{\mathcal{A}}$ for
	the corresponding rows. By hypothesis the active constraint
	gradients together with the rows of $A$ are linearly independent
	(LICQ), strict complementarity holds, so
	$\bar\lambda_j>0$ for $j\in\mathcal{A}$ and $\bar\lambda_j=0$
	otherwise, and the Hessian $Q$ is positive definite on the null
	space of the active constraints.

	\emph{Step 1: reduction to a square nonlinear system.} Under strict
	complementarity the inactive inequalities are slack at
	$\bar{\bm{z}}$ and, by continuity of
	$\bm{\mu}\mapsto(G,\bm{g})$ from Lemma~\ref{lem:data-smooth},
	remain slack for $\bm{\mu}$ in a neighbourhood of
	$\bar{\bm{\mu}}$; their multipliers remain zero. The
	Karush--Kuhn--Tucker conditions therefore reduce, locally, to the
	square system
	\begin{equation}
		\label{eq:app-kkt}
		F(\bm{z},\bm{\lambda}_{\mathcal{A}},\bm{\nu};\,\bm{\mu})
		:=
		\begin{bmatrix}
			Q\bm{z}+\bm{c}
			+G_{\mathcal{A}}(\bm{\mu})^{\top}\bm{\lambda}_{\mathcal{A}}
			+A^{\top}\bm{\nu}                                       \\[2pt]
			G_{\mathcal{A}}(\bm{\mu})\bm{z}-\bm{g}_{\mathcal{A}}(\bm{\mu}) \\[2pt]
			A\bm{z}-\bm{b}
		\end{bmatrix}
		=\bm{0},
	\end{equation}
	which the solution satisfies at $\bar{\bm{\mu}}$. By
	Lemma~\ref{lem:data-smooth}, $F$ is $C^{\infty}$ in $\bm{\mu}$ and
	affine, hence $C^{\infty}$, in
	$(\bm{z},\bm{\lambda}_{\mathcal{A}},\bm{\nu})$.

	\emph{Step 2: non-singularity of the KKT Jacobian.} The Jacobian
	of~\eqref{eq:app-kkt} in the primal--dual variables is the saddle
	matrix
	\begin{equation}
		\label{eq:app-jac}
		J=
		\begin{bmatrix}
			Q                  & G_{\mathcal{A}}^{\top} & A^{\top} \\
			G_{\mathcal{A}}    & 0                      & 0        \\
			A                  & 0                      & 0
		\end{bmatrix}.
	\end{equation}
	Let $C=\big[G_{\mathcal{A}};A\big]$ denote the stacked active and
	equality rows, which has full row rank by LICQ. Suppose
	$J(\bm{u},\bm{w})=\bm{0}$ with $\bm{u}\in\R^{n_z}$ and $\bm{w}$
	collecting the dual components. The lower blocks give
	$C\bm{u}=\bm{0}$, so $\bm{u}$ lies in the null space of $C$; the
	upper block gives $Q\bm{u}+C^{\top}\bm{w}=\bm{0}$, and premultiplying
	by $\bm{u}^{\top}$ yields
	$\bm{u}^{\top}Q\bm{u}+(C\bm{u})^{\top}\bm{w}
		=\bm{u}^{\top}Q\bm{u}=0$. Because $Q$ is positive definite on
	$\ker C$ by the reduced-Hessian hypothesis, $\bm{u}=\bm{0}$. Then
	$C^{\top}\bm{w}=\bm{0}$, and full row rank of $C$ forces
	$\bm{w}=\bm{0}$. Hence $J$ is non-singular.

	\emph{Step 3: implicit function theorem.} $F$ is continuously
	differentiable in a neighbourhood of
	$\big((\bar{\bm{z}},\bar{\bm{\lambda}}_{\mathcal{A}},\bar{\bm{\nu}}),
		\bar{\bm{\mu}}\big)$, vanishes there, and has non-singular
	Jacobian $J$ in the primal--dual variables by Step~2. The implicit
	function theorem therefore supplies a neighbourhood of
	$\bar{\bm{\mu}}$ and a continuously differentiable map
	$\bm{\mu}\mapsto
		\big(\bm{z}(\bm{\mu}),\bm{\lambda}_{\mathcal{A}}(\bm{\mu}),
		\bm{\nu}(\bm{\mu})\big)$ satisfying~\eqref{eq:app-kkt}
	identically, with
	\begin{equation}
		\label{eq:app-sens}
		\begin{bmatrix}
			\partial\bm{z}/\partial\bm{\mu} \\
			\partial\bm{\lambda}_{\mathcal{A}}/\partial\bm{\mu} \\
			\partial\bm{\nu}/\partial\bm{\mu}
		\end{bmatrix}
		=-J^{-1}\,
		\frac{\partial F}{\partial\bm{\mu}} ,
		\qquad
		\frac{\partial F}{\partial\bm{\mu}}
		=
		\begin{bmatrix}
			\big(\partial G_{\mathcal{A}}/\partial\bm{\mu}\big)^{\top}
			\bar{\bm{\lambda}}_{\mathcal{A}}                              \\
			\big(\partial G_{\mathcal{A}}/\partial\bm{\mu}\big)\bar{\bm{z}}
			-\partial\bm{g}_{\mathcal{A}}/\partial\bm{\mu}                 \\
			0
		\end{bmatrix},
	\end{equation}
	where the blocks of $\partial F/\partial\bm{\mu}$ are evaluated
	through~\eqref{eq:app-dh}. Shrinking the neighbourhood if necessary
	so that the inactive inequalities remain slack and the active
	multipliers remain strictly positive, the active set is locally
	constant and the solution of~\eqref{eq:app-qp} coincides with that
	of~\eqref{eq:app-kkt}. Since the commanded acceleration
	$\bm{a}_{t\mid t}$ is a fixed coordinate projection of
	$\bm{z}(\bm{\mu})$, the solution map~\eqref{eq:solution-map} is
	continuously differentiable at $\bar{\bm{\mu}}$, and its derivative
	is obtained by solving the linear system~\eqref{eq:app-sens}, which
	is the linearized KKT system of the quadratic program at its
	solution.
\end{proof}

\begin{remark}[Where the conclusion lapses]
	\label{rem:gradient-scope}
	The result is local and conditional, and the experiments are
	designed not to rely on more. At a change of active set the
	solution map remains continuous but is generally only directionally
	differentiable, and the derivative computed
	from~\eqref{eq:app-sens} is then a one-sided object; strict
	complementarity is what excludes this case locally. Degenerate
	solutions violating LICQ, and a reduced Hessian that is only
	semi-definite, both break Step~2. The saturations of the inner
	control law contribute their own switching surfaces, on which the
	chain from the realized trajectory to $\bm{a}_{t\mid t}$ is
	non-differentiable and, within saturated regions, has vanishing
	derivative. Training therefore relies on almost-everywhere
	differentiability rather than on global smoothness, which is
	standard for optimization layers and is the sense in which
	Proposition~\ref{prop:gradient} is used.
\end{remark}

\subsection{Support of the operational gradient}
\label{app:support}

\begin{proof}[Proof of Proposition~\ref{prop:gradient-support}]
	Let $\mathcal{J}$ denote the operational part of~\eqref{eq:tasl},
	that is, the terms
	$\lambda_s\widetilde\Phi_{\mathrm{los}}
		+\lambda_d\Phi_{\mathrm{dev}}
		+\lambda_u\Phi_{\mathrm{effort}}$ evaluated on the realized
	closed-loop trajectory, excluding the displacement anchor. Fix a
	neighbour $i$, a planned step $k$, and consider the derivative of
	$\mathcal{J}$ with respect to the predicted mean position
	$\hat{\bm{p}}^{\,i}_{k}$.

	\emph{Step 1: the forecast acts only through the solution map.}
	Under Assumption~\ref{ass:dynamics} the realized trajectory depends
	on the forecast only through the executed controls, and by
	\eqref{eq:method-composition} the control at instant $t$ is
	$\bu_t=\mathrm{Inner}(\bx_t,\bm{a}_{t\mid t})$ with
	$\bm{a}_{t\mid t}$ the first commanded acceleration
	of~\eqref{eq:solution-map}. The forecast enters no other term of
	$\mathcal{J}$: the clearance $\bar h_t$
	in~\eqref{eq:soft-los} is computed from the \emph{true} neighbour
	positions, and $\Phi_{\mathrm{dev}},\Phi_{\mathrm{effort}}$ are
	functions of the realized state and control. Hence the chain rule
	gives
	\begin{equation}
		\label{eq:app-chain}
		\frac{\partial\mathcal{J}}
		{\partial\hat{\bm{p}}^{\,i}_{k}}
		=\frac{\partial\mathcal{J}}{\partial\bm{a}_{t\mid t}}\,
		\frac{\partial\bm{a}_{t\mid t}}
		{\partial\hat{\bm{p}}^{\,i}_{k}},
	\end{equation}
	so it suffices to show that the second factor vanishes under the
	stated condition.

	\emph{Step 2: inactive rows contribute nothing.} By
	Proposition~\ref{prop:gradient} and~\eqref{eq:app-sens}, the
	sensitivity of the solution to the forecast is
	$\partial\bm{z}/\partial\bm{\mu}
		=-\big[J^{-1}\,\partial F/\partial\bm{\mu}\big]_{\bm{z}}$, and
	$\partial F/\partial\bm{\mu}$ involves only the active rows
	$G_{\mathcal{A}},\bm{g}_{\mathcal{A}}$. Consequently, if
	$\hat{\bm{p}}^{\,i}_{k}$ appears in no active row, then every block
	of $\partial F/\partial\bm{\mu}$ in the corresponding columns is
	zero: the first block carries the factor
	$\bar{\bm{\lambda}}_{\mathcal{A}}$ restricted to active rows, and
	the second and third blocks involve only
	$\partial G_{\mathcal{A}}/\partial\bm{\mu}$ and
	$\partial\bm{g}_{\mathcal{A}}/\partial\bm{\mu}$. Equivalently, in
	the full-multiplier form of the optimality system the inactive rows
	enter multiplied by $\bar\lambda_j=0$, so their parametric
	derivatives are annihilated. Hence
	$\partial\bm{a}_{t\mid t}/\partial\hat{\bm{p}}^{\,i}_{k}=\bm{0}$
	whenever no active inequality depends on $\hat{\bm{p}}^{\,i}_{k}$.

	\emph{Step 3: identifying which rows can depend on the forecast.}
	The predicted position $\hat{\bm{p}}^{\,i}_{k}$ enters
	\eqref{eq:app-qp} only through the clearance $h^{i}_{k}$, by
	\eqref{eq:cbf-barrier} and Lemma~\ref{lem:data-smooth}; the
	equality block, the acceleration bounds, the slack sign conditions,
	and the cost data are all independent of it. The clearance
	$h^{i}_{k}$ appears in the barrier
	recursion~\eqref{eq:cbf-condition} exactly twice: on the
	left-hand side of the inequality indexed $k$, and on the right-hand
	side of the inequality indexed $k+1$. For the boundary indices, as
	noted after~\eqref{eq:planning-qp}, $h^{i}_{0}$ enters only the
	inequality at $k=1$ and $h^{i}_{H_p}$ only the inequality at
	$k=H_p$. Therefore the only rows whose data depend on
	$\hat{\bm{p}}^{\,i}_{k}$ are those two, and by Step~2 the derivative
	vanishes unless at least one of them is active at the optimum.

	\emph{Step 4: the residual signal.} When the derivative in
	\eqref{eq:app-chain} vanishes, the total gradient of~\eqref{eq:tasl}
	with respect to that predicted position reduces to the contribution
	of the displacement anchor, namely
	$\lambda_a\,\partial\mathcal{L}_{\mathrm{ADE}}
		/\partial\hat{\bm{p}}^{\,i}_{k}$, since the anchor is evaluated
	directly on the forecast and does not pass through the planner.
	Composing with $\partial\hat{\bm{p}}^{\,i}_{k}/\partial\theta$ and
	summing over neighbours and steps gives the statement for
	$\nabla_\theta$.
\end{proof}

\begin{remark}[Interpretation and what it does not assert]
	\label{rem:support-scope}
	Proposition~\ref{prop:gradient-support} localizes the operational
	signal but says nothing about its usefulness. Two consequences
	matter for reading Section~\ref{sec:results-signal}. First, the
	active set is a property of the planner configuration
	$\phi=(w_{\mathrm{track}},w_{\mathrm{acc}},w_{\mathrm{slack}},
		\gamma,a_{\max},H_p)$ as much as of the encounter, so the region
	in which the objective can act at all changes when $\phi$ changes.
	This is the formal counterpart of the cross-configuration
	sensitivity reported in
	Section~\ref{sec:results-boundary}, and it is why a predictor
	trained through one planner has no structural reason to retain its
	operational effect under another. Second, the proposition constrains
	the support of the gradient, not the direction or magnitude of the
	resulting parameter change; whether accuracy is in fact reallocated
	toward the constraint-relevant region is
	Hypothesis~\ref{hyp:reallocation}, which is tested empirically and,
	in the configuration studied here, not supported. Because per-step
	active-set logs were not retained during training, the CPA-centred
	window used in
	Section~\ref{sec:results-signal} serves as a proxy for the active
	region identified here, and the mismatch between proxy and true
	support is a limitation of that analysis rather than of this
	proposition.
\end{remark}

\end{document}